%% toc: P3 36 36.8,36.9
\section{Proofs for definable second-order quantifiers and the comparison of theories}
\label{sec:annex-p03-w03-p03-046-dsoq}

Throughout, \(T\) is a first-order theory in a language \(L\) with no finite
models, and interpretations are those of
\cref{def:p03-w03-p03-046-interpretation} under the reading of
\cref{def:p03-w03-p03-046-reading}. For a binary relation variable \(s\)
interpreted by the graph of an involution, \(s(x)\) abbreviates the unique
\(w\) with \(s(x,w)\), and \(\operatorname{Mov}_s(x)\) abbreviates
\(\exists w\,(s(x,w)\wedge w\neq x)\); both are first-order in \(s\).
Cardinals are those of ZFC, and an infinite cardinal \(\kappa\) satisfies
\(\kappa+\kappa=\kappa\cdot\kappa=\kappa\).

\begin{annexitem}{annex:p03-w03-p03-046-spectrum-monotone}{thm:p03-w03-p03-046-spectrum-monotone}{Spectrum-preserving interpretations}
\begin{proof}
Let \(*\) be a spectrum-preserving interpretation of \(T_1\) into \(T_2\);
the sentences \(\phi\) range over those of the language of \(T_1\) and
\(\phi^*\) over those of the language of \(T_2\). Models are infinite, as
in \cref{def:p03-w03-p03-046-spectrum}.

\emph{L\"owenheim numbers.} Let \(\lambda=ls^{T_2}_{\mathscr L_2}\). Let
\(\phi\) be an \(\mathscr L_1\)-sentence such that \(T_1\cup\{\phi\}\) has a
model. Then \(\operatorname{Spec}_{T_1}(\phi)\neq\varnothing\), hence
\(\operatorname{Spec}_{T_2}(\phi^*)\neq\varnothing\), so
\(T_2\cup\{\phi^*\}\) has a model and therefore one of some cardinality
\(\mu\le\lambda\). Then
\(\mu\in\operatorname{Spec}_{T_2}(\phi^*)=\operatorname{Spec}_{T_1}(\phi)\),
so \(T_1\cup\{\phi\}\) has a model of cardinality at most \(\lambda\). Thus
\(\lambda\) has the property whose least instance defines
\(ls^{T_1}_{\mathscr L_1}\); the cardinals being well ordered, that least
instance exists and is at most \(\lambda\).

\emph{Hanf numbers.} Let \(\eta=h^{T_2}_{\mathscr L_2}\), and let \(\phi\)
be an \(\mathscr L_1\)-sentence such that \(T_1\cup\{\phi\}\) has a model
of power \(\eta\). Then
\(\eta\in\operatorname{Spec}_{T_1}(\phi)=\operatorname{Spec}_{T_2}(\phi^*)\),
so \(T_2\cup\{\phi^*\}\) has arbitrarily large models, that is,
\(\operatorname{Spec}_{T_2}(\phi^*)\) is unbounded. The equal set
\(\operatorname{Spec}_{T_1}(\phi)\) is then unbounded, so
\(T_1\cup\{\phi\}\) has arbitrarily large models. Hence \(\eta\) has the
defining property of \(h^{T_1}_{\mathscr L_1}\), which therefore exists and
is at most \(\eta\).
\end{proof}
\end{annexitem}

\begin{annexitem}{annex:p03-w03-p03-046-coding}{thm:p03-w03-p03-046-coding}{Coding and monadic logic}
\begin{proof}
Let \(M_0\entails T\), let the sets \(B_0\), \(B_1\) and \(C\) be infinite
subsets of \(M_0\), and let
\(\phi(x_0,x_1,y,\bar z,P_0,P_1,P_2)\) with constants \(\bar d\) define a
one-to-one function from \(B_0\times B_1\) onto \(C\), the predicates
\(P_0,P_1,P_2\) being interpreted by \(B_0,B_1,C\). Write
\(F(x,w,v)\) for \(\phi(x,w,v,\bar z,P_0,P_1,P_2)\).

\emph{Step 1: large models.} Fix \(\kappa\ge\aleph_0\). Expand \(L\) by the
predicates \(P_0,P_1,P_2\), constants for \(\bar d\), and new constants
\(e_\xi,e'_\xi\) for \(\xi<\kappa\). Let \(\Gamma\) consist of the complete
theory of \((M_0,B_0,B_1,C,\bar d)\) together with \(P_0(e_\xi)\),
\(P_1(e'_\xi)\) for \(\xi<\kappa\) and \(e_\xi\neq e_\eta\),
\(e'_\xi\neq e'_\eta\) for \(\xi\neq\eta\). A finite subset of \(\Gamma\)
mentions finitely many new constants; interpreting them by distinct members
of the infinite sets \(B_0\) and \(B_1\) makes \((M_0,B_0,B_1,C,\bar d)\) a
model of it. By \cref{thm:p03-compactness-ultraproduct} there is a model \(N\) of
\(\Gamma\). Its \(L\)-reduct is a model of \(T\), because \(T\) is
contained in the theory of \(M_0\). The first-order sentence saying that
\(F(x,w,v,\bar d)\) defines a one-to-one function from
\(P_0\times P_1\) is true in \(M_0\), hence in \(N\), and
\(|P_0^N|,|P_1^N|\ge\kappa\).

\emph{Step 2: the interpreting formulas.} All parameters for the
monadic quantifier are sets, so the parameter condition \(\psi\) is the
true formula of one unary variable. The individual parameters are
\(\bar z\); the set parameters are \(P_0,P_1,P_2\), a set \(D\) (the
domain), a set \(G\) (coding an injection of \(D\) into \(P_1\)) and a set
\(H\) (coding a non-surjective injection of \(D\) into itself). Put
\begin{align*}
\mathrm{In}(x,w,V)&:=\exists v\,(F(x,w,v)\wedge V(v)),\\
\mathrm{Gr}(x,w)&:=P_1(w)\wedge\mathrm{In}(x,w,G),\\
\mathrm{Rel}_V(x,y)&:=D(x)\wedge D(y)\wedge\exists w\,(\mathrm{Gr}(y,w)\wedge\mathrm{In}(x,w,V)).
\end{align*}
Let \(\pi\) be the conjunction of:
\begin{enumerate}
  \item \(\forall x\,(D(x)\to P_0(x))\);
  \item \(\forall x\forall w\,(D(x)\wedge P_1(w)\to\exists!v\,F(x,w,v))\), and
        for \(x,x'\in D\), \(w,w'\in P_1\), if \(F(x,w,v)\) and
        \(F(x',w',v)\) then \(x=x'\) and \(w=w'\);
  \item \(\forall x\,(D(x)\to\exists!w\,\mathrm{Gr}(x,w))\), and for
        \(x,x'\in D\), if \(\mathrm{Gr}(x,w)\) and \(\mathrm{Gr}(x',w)\) then
        \(x=x'\);
  \item \(\forall x\,(D(x)\to\exists!y\,\mathrm{Rel}_H(x,y))\), for all
        \(x,x',y\), if \(\mathrm{Rel}_H(x,y)\) and \(\mathrm{Rel}_H(x',y)\)
        then \(x=x'\), and \(\exists y\,(D(y)\wedge\neg\exists
        x\,\mathrm{Rel}_H(x,y))\).
\end{enumerate}
Let \(\theta(x):=D(x)\), and let the coding formula be
\(\chi(x_0,x_1):=\mathrm{Rel}_{V}(x_0,x_1)\), where \(V\) is a coding
parameter (a set ranging over all subsets of the model).

\emph{Clause~(1).} Let \(M\entails T\) and let parameters satisfy \(\pi\).
By conjunct~2, for \(x\in D\) and \(w\in P_1\) there is a unique \(v\) with
\(F(x,w,v)\); call it \(f(x,w)\). Then \(f\) is one-to-one on
\(D\times P_1\), and for such \(x,w\), \(\mathrm{In}(x,w,V)\) holds iff
\(f(x,w)\in V\). By conjunct~3, \(\mathrm{Gr}\) is the graph of a one-to-one
function \(g\colon D\to P_1\). Hence for \(x,y\in D\),
\(\mathrm{Rel}_V(x,y)\) holds iff \(f(x,g(y))\in V\). By conjunct~4,
\(\mathrm{Rel}_H\) is the graph of a one-to-one function \(h\colon D\to D\)
whose range omits some member of \(D\); in particular \(D\neq\varnothing\).
A one-to-one map of a finite set into itself is onto, so \(B=D\) is
infinite. Every \(\mathrm{Rel}_V\) is a binary relation on \(D\). Given
any \(R\subseteq D\times D\), put \(V=\{f(x,g(y)):(x,y)\in R\}\). For
\(x,y\in D\), \(f(x,g(y))\in V\) iff \(f(x,g(y))=f(x',g(y'))\) for some
\((x',y')\in R\), iff \((x,y)\in R\), because \(f\) and \(g\) are
one-to-one. So \(\mathrm{Rel}_V=R\), and the coded relations are exactly
the binary relations on \(B\).

\emph{Clause~(2).} Let \(\kappa\ge\aleph_0\) and take \(N\) from Step~1,
with \(P_0,P_1,P_2\) and \(\bar z\) interpreted as there. Let \(f\) be the
one-to-one function from \(P_0^N\times P_1^N\) defined by \(F\). Choose
\(D\subseteq P_0^N\) with \(|D|=\kappa\), a one-to-one
\(g\colon D\to P_1^N\) (possible since \(|P_1^N|\ge\kappa\)), and a
one-to-one \(h\colon D\to D\) that is not onto (identify \(D\) with
\(\kappa\) and let \(h(\xi)=\xi+1\), which omits \(0\)). Put
\(G=\{f(x,g(x)):x\in D\}\) and \(H=\{f(x,g(h(x))):x\in D\}\). For \(x\in D\)
and \(w\in P_1^N\), \(f(x,w)\in G\) iff \(f(x,w)=f(x',g(x'))\) for some
\(x'\in D\), iff \(w=g(x)\); so \(\mathrm{Gr}\) is the graph of \(g\) on
\(D\). For \(x,y\in D\), \(\mathrm{Rel}_H(x,y)\) iff
\(f(x,g(y))=f(x',g(h(x')))\) for some \(x'\in D\), iff \(y=h(x)\), using
that \(g\) is one-to-one; so \(\mathrm{Rel}_H\) is the graph of \(h\).
Conjunct~1 holds by choice of \(D\), conjunct~2 by Step~1, and conjuncts~3
and~4 by the properties of \(g\) and \(h\). Hence \(\pi\) holds and
\(|\theta(N)|=|D|=\kappa\).
\end{proof}
\end{annexitem}

\begin{annexitem}{annex:p03-w03-p03-046-pairing}{thm:p03-w03-p03-046-pairing}{A definable pairing function interprets second-order logic monadically}
\begin{proof}
Let \(M\) be a model of \(T\); it exists since \(T\) is consistent and is
infinite since \(T\) has no finite models. In \(M\) the formula \(\rho\)
defines a function \(\rho\colon M\times M\to M\). It is one-to-one: if
\(\rho(x,y)=\rho(x',y')\) then \(x=\rho_0(\rho(x,y))=\rho_0(\rho(x',y'))=x'\)
and likewise \(y=y'\) using \(\rho_1\). It is onto: every \(x\) equals
\(\rho(\rho_0(x),\rho_1(x))\). Hence \(T\) admits coding in the sense of
\cref{def:p03-w03-p03-046-pairing} with \(B_0=B_1=C=M\), the formula
\(\phi(x_0,x_1,y):=\rho(x_0,x_1,y)\) (in which the predicates and \(\bar z\)
do not occur) and empty \(\bar d\). \Cref{thm:p03-w03-p03-046-coding} gives
\((T_\infty,\mathrm{2nd})\le(T,\mathit{Mon})\).
\end{proof}
\end{annexitem}

\begin{annexitem}{annex:p03-w03-p03-046-equiv-uniform}{thm:p03-w03-p03-046-equiv-uniform}{Large equivalence relations and permutational logic}
\begin{proof}
All parameters for the permutational quantifier are graphs of involutions.
The individual parameters are \(\bar z\); the involution parameters are
\(\sigma\) (of the hypothesis), \(\alpha,\gamma,\delta,\varepsilon\) and
\(\zeta\); the coding parameter is an involution \(\beta\). Put
\begin{align*}
E(x,y)&:=\phi(x,y,\bar z,\sigma),\qquad
U_0(x):=\operatorname{Mov}_\gamma(x),\\
U_1(x)&:=\operatorname{Mov}_\alpha(x)\wedge\neg\operatorname{Mov}_\gamma(x),\\
\Theta(x)&:=\operatorname{Mov}_\delta(x),\\
\mathrm{Code}_\beta(u,v)&:=\Theta(u)\wedge\Theta(v)\\
&\qquad{}\wedge\exists x\,\bigl(U_0(x)\wedge\operatorname{Mov}_\beta(x)\\
&\qquad\qquad{}\wedge E(x,u)\wedge E(\alpha(x),\varepsilon(v))\bigr).
\end{align*}
Let \(\pi\) be the conjunction of:
\begin{enumerate}
  \item[(a)] \(\forall x\,(U_0(x)\to U_1(\alpha(x)))\);
  \item[(b)] \(\forall x\,(\Theta(x)\to U_0(x))\) and
        \(\forall x\,(U_0(x)\to\exists!u\,(\Theta(u)\wedge E(x,u)))\);
  \item[(c)] \(\forall u\,(\Theta(u)\to U_1(\varepsilon(u)))\) and
        \(\forall y\,\bigl(U_1(y)\to\exists!w\,(\exists u\,(\Theta(u)\wedge
        w=\varepsilon(u))\wedge E(y,w))\bigr)\);
  \item[(d)] \(\forall u\forall v\,\bigl(U_0(u)\wedge U_1(v)\to\exists!x\,(U_0(x)\wedge
        E(x,u)\wedge E(\alpha(x),v))\bigr)\);
  \item[(e)] \(\mathrm{Code}_\zeta\) is the graph of a function from
        \(\Theta\) to \(\Theta\) that is total and one-to-one, and some
        member of \(\Theta\) is not in its range.
\end{enumerate}
Let \(\theta:=\Theta\) and let the coding formula be
\(\chi(u,v):=\mathrm{Code}_\beta(u,v)\).

\emph{The coding argument.} Suppose \(M\entails T\) and parameters satisfy
(a)--(d); write \(B=\Theta(M)\), \(U_0=U_0(M)\),
\(U_1=U_1(M)\). By definition \(U_0\cap U_1=\varnothing\). Let
\(R\subseteq B\times B\), and put
\[
V=\{x\in U_0:\exists(u',v')\in R\ (E(x,u')\wedge E(\alpha(x),\varepsilon(v')))\}.
\]
For \(x\in V\), \(\alpha(x)\in U_1\) by (a), so \(\alpha(x)\neq x\), and
the two-element sets \(\{x,\alpha(x)\}\) for \(x\in V\) are pairwise
disjoint: \(\alpha\) is one-to-one, and \(\alpha(x)=x'\) with \(x,x'\in V\)
is impossible since \(\alpha(x)\in U_1\) and \(x'\in U_0\). Let \(\beta\)
exchange \(x\) and \(\alpha(x)\) for each \(x\in V\) and fix all other
points. Then \(\beta\) is an involution and
\(\operatorname{Mov}(\beta)\cap U_0=V\), since
\(\operatorname{Mov}(\beta)=V\cup\alpha[V]\) and \(\alpha[V]\subseteq U_1\).
We show \(\mathrm{Code}_\beta(M)=R\). Let \(u,v\in B\).

Suppose \((u,v)\in R\). Then \(u\in U_0\) by (b) and \(\varepsilon(v)\in U_1\) by (c),
so (d) gives \(x\in U_0\) with \(E(x,u)\) and \(E(\alpha(x),\varepsilon(v))\).
This \(x\) lies in \(V\), witnessed by \((u,v)\), so
\(\operatorname{Mov}_\beta(x)\) and \(\mathrm{Code}_\beta(u,v)\) hold.

Suppose \(\mathrm{Code}_\beta(u,v)\). Then there is \(x\in U_0\cap
\operatorname{Mov}(\beta)=V\) with \(E(x,u)\) and
\(E(\alpha(x),\varepsilon(v))\), and, as \(x\in V\), some \((u',v')\in R\)
with \(E(x,u')\) and \(E(\alpha(x),\varepsilon(v'))\). Both \(u\) and \(u'\)
lie in \(B\) and are \(E\)-related to \(x\), so \(u=u'\) by the uniqueness
in (b). The point \(\alpha(x)\) lies in \(U_1\) by (a), and both
\(\varepsilon(v)\) and \(\varepsilon(v')\) lie in \(\varepsilon[B]\) and are
\(E\)-related to it, so \(\varepsilon(v)=\varepsilon(v')\) by the
uniqueness in (c), and \(v=v'\) since \(\varepsilon\) is a bijection. Thus
\((u,v)\in R\).

\emph{Clause~(1).} Let \(M\entails T\) and let the parameters satisfy
\(\pi\). By (e), \(\mathrm{Code}_\zeta\) is the graph of a one-to-one map
of \(B\) into \(B\) that is not onto; so \(B\neq\varnothing\), and \(B\) is
infinite since a one-to-one self-map of a finite set is onto. Each
\(\mathrm{Code}_\beta(M)\) is a binary relation on \(B\) by its first two
conjuncts, and by the coding argument every binary relation on \(B\) is
some \(\mathrm{Code}_\beta(M)\).

\emph{Clause~(2).} Let \(\kappa\ge\aleph_0\), and take \(M\), \(\bar c\),
\(\sigma\) and \(A\) as in the hypothesis; interpret \(\bar z\) by
\(\bar c\). For \(x,y\in A\), \(E(x,y)\) holds iff \(x\) and \(y\) lie in
the same class. Since there are at least \(\kappa=\kappa+\kappa\) classes,
choose pairwise distinct classes \(K_{0,\xi}\), \(K_{1,\eta}\)
(\(\xi,\eta<\kappa\)), and in them pairwise distinct points
\(x_{\xi,\eta}\in K_{0,\xi}\) and \(y_{\eta,\xi}\in K_{1,\eta}\)
(\(\xi,\eta<\kappa\)), possible as each class has at least \(\kappa\)
members. Put \(U_0=\{x_{\xi,\eta}\}\), \(U_1=\{y_{\eta,\xi}\}\); these are
disjoint and of cardinality \(\kappa\cdot\kappa=\kappa\). Define the
involutions:
\(\alpha\) exchanges \(x_{\xi,\eta}\) and \(y_{\eta,\xi}\) for all
\(\xi,\eta\) and fixes all other points;
\(\gamma\) has \(\operatorname{Mov}(\gamma)=U_0\), obtained by splitting
\(U_0\) into two-element blocks (\(|U_0|=|U_0\times 2|\)) and exchanging
the members of each block;
\(\delta\) likewise has \(\operatorname{Mov}(\delta)=B:=\{x_{\xi,0}:\xi<\kappa\}\);
\(\varepsilon\) exchanges \(x_{\xi,0}\) and \(y_{\xi,0}\) for each \(\xi\)
and fixes all other points. Then \(U_0(M)=U_0\),
\(U_1(M)=\operatorname{Mov}(\alpha)\setminus\operatorname{Mov}(\gamma)=U_1\)
and \(\Theta(M)=B\). We check (a)--(d).
(a) \(\alpha(x_{\xi,\eta})=y_{\eta,\xi}\in U_1\).
(b) \(B\subseteq U_0\); for \(x_{\xi,\eta}\in U_0\) the members of \(B\) in
its class \(K_{0,\xi}\) are exactly \(x_{\xi,0}\), since the classes
\(K_{0,\xi'}\) are distinct.
(c) \(\varepsilon(x_{\xi,0})=y_{\xi,0}\in U_1\); for \(y_{\eta,\xi}\in U_1\)
the members of \(\varepsilon[B]=\{y_{\xi',0}\}\) in its class
\(K_{1,\eta}\) are exactly \(y_{\eta,0}\).
(d) If \(u\in K_{0,\xi}\) and \(v\in K_{1,\eta}\), the points
\(x\in U_0\) with \(E(x,u)\) are the \(x_{\xi,\eta'}\), and
\(\alpha(x_{\xi,\eta'})=y_{\eta',\xi}\) is \(E\)-related to \(v\) iff
\(\eta'=\eta\); so \(x_{\xi,\eta}\) is the unique such point.
Finally let \(h\colon B\to B\) be one-to-one and not onto (identify \(B\)
with \(\kappa\) and use \(\xi\mapsto\xi+1\)), and let \(\zeta\) be the
involution that the coding argument produces for the graph of \(h\); that
argument used only (a)--(d), so \(\mathrm{Code}_\zeta(M)\) is
the graph of \(h\) and (e) holds. Hence \(\pi\) holds and
\(|\theta(M)|=|B|=\kappa\).
\end{proof}
\end{annexitem}

\begin{annexitem}{annex:p03-w03-p03-046-equiv}{thm:p03-w03-p03-046-equiv}{An equivalence relation with infinitely many infinite classes}
\begin{proof}
Let \(M_0\entails T\), \(A_0\subseteq M_0\) and \(\phi(x,y)\) be as in the
hypothesis. Fix \(\kappa\ge\aleph_0\). Expand \(L\) by a unary predicate
\(P\) and new constants \(e_{\xi,\eta}\) (\(\xi,\eta<\kappa\)). Let
\(\Gamma\) be the complete theory of \((M_0,A_0)\) together with
\(P(e_{\xi,\eta})\), \(\phi(e_{\xi,\eta},e_{\xi,0})\),
\(\neg\phi(e_{\xi,0},e_{\xi',0})\) for \(\xi\neq\xi'\), and
\(e_{\xi,\eta}\neq e_{\xi,\eta'}\) for \(\eta\neq\eta'\). A finite subset of
\(\Gamma\) mentions constants with finitely many first indices
\(\xi_1,\dots,\xi_p\); since \(\phi\) has infinitely many infinite classes
on \(A_0\), choose \(p\) distinct infinite classes and send the constants
with first index \(\xi_i\) to distinct members of the \(i\)th class. This
satisfies the finite subset in \((M_0,A_0)\). By
\cref{thm:p03-compactness-ultraproduct}, \(\Gamma\) has a model \((N,A)\). Its
\(L\)-reduct is a model of \(T\), and the first-order sentence saying that
\(\phi\) is an equivalence relation on \(P\) holds in \(M_0\), hence in
\((N,A)\). The classes of \(e_{\xi,0}\) (\(\xi<\kappa\)) are pairwise
distinct, and each contains the \(\kappa\) distinct points
\(e_{\xi,\eta}\). So the hypothesis of
\cref{thm:p03-w03-p03-046-equiv-uniform} holds at \(\kappa\) for the formula
\(\phi(x,y)\), regarded as a formula in which \(\bar z\) and \(s\) do not
occur, with \(\sigma\) the identity involution. As \(\kappa\) was arbitrary,
\cref{thm:p03-w03-p03-046-equiv-uniform} gives
\((T_\infty,\mathrm{2nd})\le(T,\text{1-1})\).
\end{proof}
\end{annexitem}

\begin{annexitem}{annex:p03-w03-p03-046-order}{thm:p03-w03-p03-046-order}{Second-order logic in the permutational theory of order}
\begin{proof}
\(\mathit{Th}(<)\), the theory of infinite linear orders, has no finite
models. Let
\[
\phi(x,y,s):=\neg\exists z\,\bigl(\operatorname{Mov}_s(z)\wedge((x<z\wedge z\le y)\vee(y<z\wedge z\le x))\bigr).
\]
Fix \(\kappa\ge\aleph_0\). Let \(M=\kappa\times\kappa\) ordered
lexicographically: \((\xi,\eta)<(\xi',\eta')\) iff \(\xi<\xi'\), or
\(\xi=\xi'\) and \(\eta<\eta'\). This is an infinite linear order, so
\(M\entails\mathit{Th}(<)\). Let \(P=\{(\xi,0):\xi<\kappa\}\), and let
\(\sigma\) be an involution with \(\operatorname{Mov}(\sigma)=P\) (split the
infinite set \(P\) into two-element blocks). We claim that for
\(x=(\xi,\eta)\) and \(y=(\xi',\eta')\), \(\phi(x,y,\sigma)\) holds iff
\(\xi=\xi'\). If \(\xi=\xi'\), say \(\eta\le\eta'\), every \(z\) with
\(x<z\le y\) has the form \((\xi,\iota)\) with \(\iota>\eta\ge 0\), so
\(z\notin P\); the case \(\eta'\le\eta\) is the same with \(x,y\)
exchanged. If \(\xi<\xi'\), then \(z=(\xi',0)\in P\) satisfies
\(x<z\le y\), so \(\phi(x,y,\sigma)\) fails; the case \(\xi'<\xi\) is
symmetric. Hence \(\phi(x,y,\sigma)\) defines on \(A=M\) the equivalence
relation of having the same first coordinate, whose \(\kappa\) classes
\(\{\xi\}\times\kappa\) each have \(\kappa\) members. The hypothesis of
\cref{thm:p03-w03-p03-046-equiv-uniform} therefore holds at every
\(\kappa\ge\aleph_0\), with empty \(\bar c\), and that lemma gives
\((T_\infty,\mathrm{2nd})\le(\mathit{Th}(<),\text{1-1})\).
\end{proof}
\end{annexitem}

\subsection*{Proofs for definable second-order quantifiers over pure sets}

The following proofs concern \cref{sec:p03-w03-p03-046-baldwin}.
Interpretations are those of \cref{def:p03-w03-p03-046-b-interpretation}
under \cref{def:p03-w03-p03-046-b-reading}. Two facts are used
repeatedly.
\begin{enumerate}
  \item[(F1)] \emph{Invariance.} If \(\psi\) is a formula of pure identity
        theory and \(g\colon A\to A'\) is a bijection, then
        \(A\entails\psi(\bar R)\) iff \(A'\entails\psi(g[\bar R])\); more
        generally \(g\) is an isomorphism from \((A,\bar b,\bar R)\) onto
        \((A',g\bar b,g[\bar R])\), so it preserves every first-order
        formula and every formula with definable second-order quantifiers,
        whose semantics refer only to the underlying set.
  \item[(F2)] \emph{Substitution.} If \(\theta(\bar x,\bar w)\) is a
        formula and \(\sigma(r)\) a formula with the relation variable
        \(r\) free, then \(\sigma[\theta/r]\), obtained by replacing each
        atom \(r(\bar t)\) by \(\theta(\bar t,\bar w)\) after renaming bound
        variables of \(\sigma\) away from \(\bar w\), holds under an
        assignment iff \(\sigma\) holds when \(r\) is assigned the relation
        that \(\theta\) defines with the same values of \(\bar w\). This is
        proved by induction on \(\sigma\); every quantifier clause, first-
        or second-order, refers to the assignment only through the values of
        the variables, so each step is immediate from the induction
        hypothesis.
\end{enumerate}
Whenever a guard \(\chi\) has relation variables \(\bar s_0\), it contains
the conjunct \(\phi(s)\) for each \(s\) among them, where \(Q_\phi\) is the
interpreting quantifier; this makes \(\bar S_0\in\mathscr R_\phi(A)\) part
of what \(A\entails\chi(\bar a_0,\bar S_0)\) says, as clause~(i) requires.
This conjunct is not repeated below.

\begin{annexitem}{annex:p03-w03-p03-046-b-transfer}{thm:p03-w03-p03-046-b-transfer}{Transfer along a full-domain interpretation}
\begin{proof}
Let \((\chi,\theta_0,\theta)\) interpret \(Q_\psi\) in \(Q_\phi\), with
\(\theta_0\) the formula \(x_0=x_0\), guard variables
\(\bar y_0,\bar s_0\) and coding variables \(\bar y,\bar s\).

\emph{Guards exist in every infinite set.} Let \(M\) be infinite. Clause~(ii)
for \(B=M\) gives \(A,\bar a_0,\bar S_0\) with
\(A\entails\chi(\bar a_0,\bar S_0)\) and \(A=\theta_0(A,\bar a_0,\bar
S_0)\approx M\). For a bijection \(h\colon A\to M\), (F1) gives
\(M\entails\chi(h\bar a_0,h[\bar S_0])\).

\emph{The translation.} Define \(\sigma'\) by induction on \(\sigma\):
atoms are unchanged, \('\) commutes with connectives, first-order
quantifiers and the quantifiers of \(\mathcal Q\), and
\[
(Q_\psi r\,\sigma_0)'\;:=\;\exists\bar y_0\,Q_\phi\bar s_0\,\exists\bar
y\,Q_\phi\bar s\,\bigl(\chi(\bar y_0,\bar s_0)\wedge
\sigma_0'[\theta(\cdot,\bar y_0,\bar s_0,\bar y,\bar s)/r]\bigr),
\]
where \(Q_\phi\bar s\) abbreviates one \(Q_\phi\) for each variable of
\(\bar s\) and the new variables are chosen fresh. The free variables are
unchanged. For a formula of \(L_{\kappa\lambda}\) the translation also
commutes with conjunctions and disjunctions of fewer than \(\kappa\)
formulas, \((\bigwedge\Phi)':=\bigwedge\{\sigma':\sigma\in\Phi\}\), and with
blocks of fewer than \(\lambda\) first-order quantifiers; the clause for
\(Q_\psi\) adds finitely many quantifiers, so \(\sigma'\) lies in
\(L_{\kappa\lambda}(\mathcal Q)\).

\emph{Correctness.} By induction on \(\sigma\), for every infinite
\(L\)-structure \(M\) and assignment \(g\) (to individual variables and
to relation variables, of the right arities),
\(M\entails\sigma[g]\iff M\entails\sigma'[g]\). Only the \(Q_\psi\) step
needs comment. Suppose \(M\entails Q_\psi r\,\sigma_0[g]\): there is
\(R\in\mathscr R_\psi(M)\) with \(M\entails\sigma_0[g,r\mapsto R]\). Choose
guard parameters \(\bar a_0,\bar S_0\) in \(M\) as above. By clause~(i),
as \(\theta_0(M,\bar a_0,\bar S_0)=M\), there are \(\bar a\in M\) and
\(\bar S\in\mathscr R_\phi(M)\) with \(\theta(M,\bar a_0,\bar S_0,\bar
a,\bar S)=R\). By the induction hypothesis
\(M\entails\sigma_0'[g,r\mapsto R]\), and by (F2) the formula
\(\sigma_0'[\theta/r]\) holds at these parameters; so
\(M\entails(Q_\psi r\,\sigma_0)'[g]\). Conversely, if the translation holds,
witnessed by \(\bar a_0,\bar S_0,\bar a,\bar S\) with the guard true, then
by clause~(i) the relation \(R=\theta(M,\bar a_0,\bar S_0,\bar a,\bar S)\)
lies in \(\mathscr R_\psi(M)\); by (F2) \(M\entails\sigma_0'[g,r\mapsto
R]\), by the induction hypothesis \(M\entails\sigma_0[g,r\mapsto R]\), and
so \(M\entails Q_\psi r\,\sigma_0[g]\). The clauses for infinitary
conjunctions, disjunctions and quantifier blocks follow from the induction
hypothesis for their components. This proves part~(1).

\emph{Part~(2).} The extended notion of Definition~1.2.4 is read as
follows: the guard, domain and coding formulas are as before, there is one
interpreting formula for each relation symbol of \(L(T)\); clause~(i)
requires, when a model \(M\) of \(T\) satisfies the guard, that the domain
with the interpreted relations be a model of \(T\) and that the coded
family be \(\mathscr R_\psi\) of the domain; clause~(ii) requires every
model of \(T\) to be isomorphic to such an interpreted model
\cite[p.~448]{Baldwin1985DSOQ}. With the domain \(x_0=x_0\) and each
relation symbol interpreted by itself, the interpreted model is \(M\)
itself. If \(M\entails T\) and \(M\entails\chi(\bar a_0,\bar S_0)\), the
guard being a formula of pure identity theory with relation variables, the
set \(M\) satisfies it, and clause~(i) of the given interpretation, applied
to the set \(M\), makes the coded family \(\mathscr R_\psi(M)\). For
clause~(ii), a model \(B\) of \(T\) is infinite, so by the first paragraph
of this proof the set \(B\) has guard parameters, and \(B\) is the
interpreted model at those parameters.
\end{proof}
\end{annexitem}

\begin{annexitem}{annex:p03-w03-p03-046-b-relativize}{thm:p03-w03-p03-046-b-relativize}{Relativized transfer}
\begin{proof}
Define \(\sigma^{\mathrm{rel}}(\bar y_0,\bar s_0)\) by induction: atoms
are unchanged, \(\mathrm{rel}\) commutes with connectives,
\((\exists x\,\sigma_0)^{\mathrm{rel}}:=\exists
x\,(\theta_0(x,\bar y_0,\bar s_0)\wedge\sigma_0^{\mathrm{rel}})\), and
\[
(Q_\psi r\,\sigma_0)^{\mathrm{rel}}:=\exists\bar y\,Q_\phi\bar
s\,\sigma_0^{\mathrm{rel}}[\theta(\cdot,\bar y_0,\bar s_0,\bar y,\bar s)/r].
\]
Let \(A\entails\chi(\bar a_0,\bar S_0)\) and
\(B=\theta_0(A,\bar a_0,\bar S_0)\). By induction on \(\sigma\): for every
assignment \(g\) sending individual variables into \(B\) and relation
variables to relations on \(B\),
\(A\entails\sigma^{\mathrm{rel}}[\bar a_0,\bar S_0,g]\iff
B\entails\sigma[g]\). Atoms \(x=y\) and \(r(\bar x)\) have the same values
in \(A\) and \(B\) under \(g\). The first-order step holds because the
relativized quantifier ranges over \(B\). For the \(Q_\psi\) step,
\(B\entails Q_\psi r\,\sigma_0[g]\) iff some \(R\in\mathscr R_\psi(B)\)
satisfies \(B\entails\sigma_0[g,r\mapsto R]\); by clause~(i) these \(R\)
are exactly the relations \(\theta(A,\bar a_0,\bar S_0,\bar a,\bar S)\)
with \(\bar a\in A\) and \(\bar S\in\mathscr R_\phi(A)\), all of which are
relations on \(B\); so by the induction hypothesis and (F2) the condition
is equivalent to \(A\entails(Q_\psi r\,\sigma_0)^{\mathrm{rel}}[\bar
a_0,\bar S_0,g]\).

Put \(\sigma^\sharp:=\exists\bar y_0\,Q_\phi\bar
s_0\,(\chi(\bar y_0,\bar s_0)\wedge\sigma^{\mathrm{rel}}(\bar y_0,\bar
s_0))\). The first claim is the displayed equivalence for sentences. If
\(B'\) is an infinite model of \(\sigma\), clause~(ii) gives \(A\),
\(\bar a_0\), \(\bar S_0\) with the guard true and
\(\theta_0(A,\bar a_0,\bar S_0)\approx B'\); by (F1) the set
\(\theta_0(A,\bar a_0,\bar S_0)\) is a model of \(\sigma\), so
\(A\entails\sigma^\sharp\), and \(A\) is infinite since it contains an
infinite set.
\end{proof}
\end{annexitem}

The next proofs use one more fact.
\begin{enumerate}
  \item[(F3)] \emph{Spreading.} If a formula \(\psi(\bar r)\) of pure identity
        theory holds of relations on some infinite set, it holds of relations
        on every infinite set. For let \(B_0\) be infinite with
        \(B_0\entails\psi(\bar R_0)\), and let \(\kappa\) be the cardinality
        of a given infinite set. The first-order theory in the
        vocabulary \(\bar R\) consisting of \(\psi(\bar R)\) and the sentences
        \(\exists^{\ge n}x\,(x=x)\), \(n<\omega\), together with
        \(\kappa\) new constants declared pairwise distinct, is finitely
        satisfiable in \((B_0,\bar R_0)\); by compactness
        (\cref{thm:p03-compactness-ultraproduct}) it has a model of
        cardinality at least \(\kappa\), and by the downward
        L\"owenheim--Skolem theorem
        (\cref{thm:p03-w03-p03-039-downward-ls}) one of cardinality exactly
        \(\kappa\). Transporting it along a bijection onto a given set of
        cardinality \(\kappa\) gives, by (F1), relations satisfying \(\psi\)
        on that set.
\end{enumerate}

\begin{annexitem}{annex:p03-w03-p03-046-b-normal}{thm:p03-w03-p03-046-b-normal}{Normalizing the coded structures}
\begin{proof}
Write \(D(x)\) for \(\theta_0(x,\bar y_0,\bar s_0)\) and, for a tuple
\(\bar p=(\bar y,\bar s)\) of coding variables,
\[
\rho(\bar x,\bar p):=\theta(\bar x,\bar y_0,\bar s_0,\bar p)\wedge
\textstyle\bigwedge_{i<k}D(x_i).
\]
Let \(\Psi(\bar p)\) be the relativization of \(\psi\) to \(D\) with \(r\)
replaced by \(\rho(\cdot,\bar p)\): every quantifier \(\exists x\,\beta\) of
\(\psi\) becomes \(\exists x\,(D(x)\wedge\beta)\) and every atom
\(r(\bar t)\) becomes \(\rho(\bar t,\bar p)\). By induction on \(\psi\), for
any set \(A\) and values of the parameters with \(D(A)\neq\varnothing\),
\(A\entails\Psi(\bar p)\) iff \((D(A),\rho(A,\bar p))\entails\psi\): the
relativized quantifiers range over \(D(A)\), and \(\rho(A,\bar p)\) is a
relation on \(D(A)\). Take a fresh copy \(\bar p_1=(\bar y_1,\bar s_1)\) of
the coding variables and put
\begin{align*}
\theta_0'(x,\bar y_0\bar y_1,\bar s_0\bar s_1)&:=\theta_0(x,\bar y_0,\bar
s_0)\wedge\Psi(\bar p_1),\\
\chi'(\bar y_0\bar y_1,\bar s_0\bar s_1)&:=\chi(\bar y_0,\bar
s_0)\wedge\Psi(\bar p_1),\\
\theta'(\bar x,\bar y_0\bar y_1,\bar s_0\bar s_1,\bar p)&:=
\bigl(\Psi(\bar p)\wedge\rho(\bar x,\bar p)\bigr)\vee
\bigl(\neg\Psi(\bar p)\wedge\rho(\bar x,\bar p_1)\bigr),
\end{align*}
with the conjunct \(\phi(s)\) in \(\chi'\) for each variable \(s\) of
\(\bar s_1\), as for every guard.

\emph{Every coded structure satisfies \(\psi\).} Fix values of all
parameters in a set \(A\) with \(\theta_0'(A,\dots)\neq\varnothing\). Then
\(A\entails\Psi(\bar p_1)\), so \(\theta_0'(A,\dots)=D(A)\neq\varnothing\)
and \((D(A),\rho(A,\bar p_1))\entails\psi\). For a coding tuple \(\bar p\),
if \(A\entails\Psi(\bar p)\) then \(\theta'(A,\dots,\bar p)=\rho(A,\bar p)\)
and \((D(A),\rho(A,\bar p))\entails\psi\); otherwise
\(\theta'(A,\dots,\bar p)=\rho(A,\bar p_1)\), which also gives a model of
\(\psi\).

\emph{Clause~(i).} Let \(A\entails\chi'(\bar a_0\bar a_1,\bar S_0\bar
S_1)\). Then \(A\entails\chi(\bar a_0,\bar S_0)\), so \(D=D(A)\) is infinite
and \(\bar S_0\in\mathscr R_\phi(A)\); also \(\bar S_1\in\mathscr
R_\phi(A)\), and \(\theta_0'(A,\dots)=D\) because \(\Psi(\bar a_1,\bar S_1)\)
holds. By the previous paragraph every coded relation lies in
\(\mathscr R_\psi(D)\). Conversely, if \(R\in\mathscr R_\psi(D)\), the
hypothesis gives \(\bar a\in A\) and \(\bar S\in\mathscr R_\phi(A)\) with
\(R=\rho(A,\bar a,\bar S)\); then \(A\entails\Psi(\bar a,\bar S)\) and
\(\theta'(A,\dots,\bar a,\bar S)=R\). So the coded family is exactly
\(\mathscr R_\psi(D)\).

\emph{Clause~(ii).} Let \(B\) be infinite. Clause~(ii) for the given
formulas yields \(A,\bar a_0,\bar S_0\) with \(A\entails\chi(\bar a_0,\bar
S_0)\) and \(D=D(A)\approx B\). By (F3), \(\mathscr
R_\psi(D)\neq\varnothing\); for \(R\) in it, the hypothesis gives \(\bar
a_1\in A\) and \(\bar S_1\in\mathscr R_\phi(A)\) with \(R=\rho(A,\bar
a_1,\bar S_1)\), so \(A\entails\Psi(\bar a_1,\bar S_1)\). Then \(A\entails
\chi'(\bar a_0\bar a_1,\bar S_0\bar S_1)\) and \(\theta_0'(A,\dots)=D\approx
B\).
\end{proof}
\end{annexitem}

\begin{annexitem}{annex:p03-w03-p03-046-b-qi-normal}{thm:p03-w03-p03-046-b-qi-normal}{Interpretations in \(Q_{\mathrm I}\) are full-domain}
\begin{proof}
Let \((\chi,\theta_0,\theta)\) interpret \(Q_\psi\) in \(Q_{\mathrm I}\).
The relations of \(Q_{\mathrm I}\) are the singletons; replacing each
relation variable \(s\) of the formulas by a new individual variable
\(z_s\), each atom \(s(t)\) by \(t=z_s\), and deleting the guard conjunct
\(\psi_{\mathrm I}(s)\), gives formulas \(\chi(\bar y_0)\),
\(\theta_0(x,\bar y_0)\), \(\theta(\bar x,\bar y_0,\bar y)\) whose values at
a point \(c\) are those of the old formulas at \(\{c\}\). So all parameters
may be taken to be points. Let \(p\) be the length of \(\bar y_0\bar y\).

\emph{Patterns.} For a set \(A\) and \(\bar q\in A^p\), two tuples
\(\bar x,\bar x'\in A^k\) have the same \emph{pattern over \(\bar q\)} if,
for all \(i,j<k\) and \(m<p\), \(x_i=x_j\iff x'_i=x'_j\) and \(x_i=q_m\iff
x'_i=q_m\). A pattern \(\pi\) is described by the conjunction
\(\varepsilon_\pi(\bar x,\bar z)\) of the corresponding equalities and
inequalities among \(x_0,\dots,x_{k-1},z_0,\dots,z_{p-1}\) that involve
some \(x_i\); there are finitely many patterns for given \(k\) and \(p\). If
\(A\) is infinite and \(\bar x,\bar x'\) have the same pattern over \(\bar
q\), there is a permutation of \(A\) fixing \(\bar q\) pointwise and
sending \(\bar x\) to \(\bar x'\): the map \(x_i\mapsto x'_i\) is well
defined and injective by the pattern, sends the points outside \(\bar q\)
to points outside \(\bar q\) and fixes the points of \(\bar q\) among the
\(x_i\); as a finite partial bijection of the infinite set \(A-\bar q\) it
extends to a permutation of \(A-\bar q\), and the identity on \(\bar q\)
completes it. Hence a set \(E\subseteq A^k\) invariant under all
permutations fixing \(\bar q\) pointwise is
\(\bigcup_{\pi\in\Sigma}\varepsilon_\pi(A,\bar q)\) for a set \(\Sigma\) of
patterns.

\emph{Part~(1).} For a set \(\Sigma\) of patterns and \(I\subseteq p\), let
\(\delta_{\Sigma,I}(\bar x,\bar z)\) be
\(\bigvee_{\pi\in\Sigma}\varepsilon^I_\pi\), where \(\varepsilon^I_\pi\)
arises from \(\varepsilon_\pi\) by replacing each atom \(x_i=z_m\) with
\(m\notin I\) by a false sentence (so that its negation becomes true). This
is a finite list. Let \(M\) be infinite and \(R\in\mathscr R_\psi(M)\).
Clause~(ii) with \(B=M\) gives \(A\) and \(\bar a_0\) with
\(A\entails\chi(\bar a_0)\) and a bijection \(g\colon M\to D\), where
\(D=\theta_0(A,\bar a_0)\). By (F1), \(g[R]\in\mathscr R_\psi(D)\), so by
clause~(i) \(g[R]=\theta(A,\bar a_0,\bar a)\) for some \(\bar a\in A\). Put
\(\bar q=\bar a_0\bar a\). The set \(E=g[R]\) is defined from \(\bar q\), so
by (F1) it is invariant under permutations fixing \(\bar q\), and
\(E=\bigcup_{\pi\in\Sigma}\varepsilon_\pi(A,\bar q)\) for some \(\Sigma\).
Let \(I=\{m<p:q_m\in D\}\), \(c_m=g^{-1}(q_m)\) for \(m\in I\), and \(c_m\)
an arbitrary point of \(M\) for \(m\notin I\). For \(\bar x\in M^k\):
\(gx_i=gx_j\iff x_i=x_j\) as \(g\) is injective; \(gx_i=q_m\iff x_i=c_m\)
for \(m\in I\); and \(gx_i\neq q_m\) for \(m\notin I\), since \(gx_i\in D\).
So \(M\entails\varepsilon^I_\pi(\bar x,\bar c)\iff
A\entails\varepsilon_\pi(g\bar x,\bar q)\), and \(\bar x\in R\iff g\bar x\in
E\iff M\entails\delta_{\Sigma,I}(\bar x,\bar c)\).

\emph{Part~(2).} List the formulas of part~(1) as
\(\delta_1,\dots,\delta_N\) and put
\[
\delta^*(\bar x,\bar z,u,\bar v):=\bigvee_{i=1}^{N}\Bigl(v_i=u\wedge
\textstyle\bigwedge_{j<i}v_j\neq u\wedge\delta_i(\bar x,\bar z)\Bigr)\vee
\Bigl(\bigwedge_{j\le N}v_j\neq u\wedge\delta_N(\bar x,\bar z)\Bigr).
\]
In an infinite set every \(\delta_i(M,\bar c)\) is a value of \(\delta^*\)
(take \(v_i=u\) and \(v_j\neq u\) for \(j<i\)), and every value of
\(\delta^*\) is some \(\delta_i(M,\bar c)\). Let
\(\Psi(\bar z,u,\bar v):=\psi[\delta^*(\cdot,\bar z,u,\bar v)/r]\). With a
fresh copy \((\bar z',u',\bar v')\) of the parameters define
\begin{align*}
\chi^*(\bar z',u',\bar v')&:=\Psi(\bar z',u',\bar v'),\qquad
\theta_0^*(x_0):=x_0=x_0,\\
\theta^*(\bar x,\bar z'u'\bar v',\bar zu\bar v)&:=\bigl(\Psi(\bar z,u,\bar
v)\wedge\delta^*(\bar x,\bar z,u,\bar v)\bigr)\vee\bigl(\neg\Psi(\bar z,u,\bar
v)\wedge\delta^*(\bar x,\bar z',u',\bar v')\bigr).
\end{align*}
\emph{Clause~(i).} If \(M\entails\chi^*(\bar e')\), the domain is \(M\),
which is infinite. By (F2) each value of \(\theta^*\) is a value of
\(\delta^*\) satisfying \(\psi\), either through the test \(\Psi\) or through
the guard; so it lies in \(\mathscr R_\psi(M)\). Every \(R\in\mathscr
R_\psi(M)\) is \(\delta_i(M,\bar c)\) by part~(1), hence a value of
\(\delta^*\) at parameters where \(\Psi\) holds, and there \(\theta^*\)
takes the value \(R\). \emph{Clause~(ii).} The coded family of the given
interpretation at the guard parameters supplied by clause~(ii) is nonempty
and equals \(\mathscr R_\psi\) of an infinite set, so by (F3)
\(\mathscr R_\psi(B)\neq\varnothing\) for every infinite \(B\); a member is
a value of \(\delta^*\) at parameters satisfying \(\chi^*\), and the domain
is \(B\) itself. The interpretation uses no relation parameters, which is
allowed since the interpreting formulas may contain any finite sequence of
relation variables, including the empty one.
\end{proof}
\end{annexitem}

\begin{annexitem}{annex:p03-w03-p03-046-b-hanf}{thm:p03-w03-p03-046-b-hanf}{Interpretations preserve Hanf numbers}
\begin{proof}
If the interpretation is in \(Q_{\mathrm I}\), replace it by a full-domain
one using \cref{thm:p03-w03-p03-046-b-qi-normal}(2). Apply
\cref{thm:p03-w03-p03-046-b-transfer}(1) with the empty vocabulary,
\(\kappa=\lambda=\omega\) and \(\mathcal Q=\{Q_\phi\}\): each
\(Q_\psi\)-sentence \(\sigma\) of pure identity theory has a
\(Q_\phi\)-sentence \(\sigma'\) true in exactly the same infinite sets. Let
\(\mu=h(Q_\phi)\). If \(\sigma\) has a model of cardinality at least
\(\mu\), so does \(\sigma'\); hence \(\sigma'\) has arbitrarily large
models, and these are models of \(\sigma\). So \(h(Q_\psi)\le\mu\).
Similarly, if \(\sigma\) has an infinite model, so does \(\sigma'\), which
then has one of cardinality at most \(\ell(Q_\phi)\), a model of
\(\sigma\); so \(\ell(Q_\psi)\le\ell(Q_\phi)\).
\end{proof}
\end{annexitem}

\begin{annexitem}{annex:p03-w03-p03-046-b-full-domain}{thm:p03-w03-p03-046-b-full-domain}{Full domains when \(Q_{1-1}\) is interpreted with full domain}
\begin{proof}
\emph{Part~(1).} Let \((\chi,\theta_0,\theta)\) interpret \(Q_\psi\) in
\(Q_\phi\), let \(A\) be infinite and \(\kappa=|A|\). Clause~(ii) for a set
of cardinality \(\kappa\) gives \(A',\bar a_0,\bar S_0\) with
\(A'\entails\chi(\bar a_0,\bar S_0)\) and \(|D'|=\kappa\), where
\(D'=\theta_0(A',\bar a_0,\bar S_0)\). Regard \(\mathfrak A'=(A',\bar
a_0,\bar S_0)\) as a first-order structure in the finite vocabulary of
constants for \(\bar a_0\) and relation symbols for \(\bar S_0\). By the
downward L\"owenheim--Skolem theorem
(\cref{thm:p03-w03-p03-039-downward-ls}, for the fragment of first-order
formulas) there is \(\mathfrak A''\preccurlyeq\mathfrak A'\) with universe
\(A''\supseteq D'\cup\{\bar a_0\}\) and \(|A''|=\kappa\). The guard is a
first-order formula of this vocabulary and contains the conjunct \(\phi(s)\)
for each of its relation variables, so by elementarity
\(\mathfrak A''\entails\chi(\bar a_0,\bar S_0\restriction A'')\), and
\(\theta_0(A'',\bar a_0,\bar S_0\restriction A'')=D'\cap A''=D'\). A
bijection \(h\colon A''\to A\) carries this, by (F1), to guard parameters
\(h\bar a_0,h[\bar S_0\restriction A'']\) in \(A\) whose domain
\(h[D']\) has cardinality \(\kappa\). No property of \(Q_\phi\) was used.

\emph{Composition.} If \(Q_a\le Q_b\) by a full-domain interpretation
\((\chi_1,x_0=x_0,\theta_1)\) and \(Q_b\le Q_c\) by a full-domain
interpretation \((\chi_2,x_0=x_0,\theta_2)\), then \(Q_a\le Q_c\) by a
full-domain interpretation. Fix a guard tuple \(\bar w_0\) for
\(\chi_2\) and, for each relation variable \(s\) of \(\chi_1\) or
\(\theta_1\), a fresh coding tuple \(\bar w_s\) (points and \(Q_c\)-relation
variables). Replace each atom \(s(\bar t)\) by
\(\theta_2(\bar t,\bar w_0,\bar w_s)\), obtaining \(\chi_1^\circ\),
\(\theta_1^\circ\), and put \(\chi:=\chi_2(\bar w_0)\wedge\chi_1^\circ\),
\(\theta:=\theta_1^\circ\). If \(\chi\) holds in \(A\), the values
\(\theta_2(A,\bar w_0,\bar w_s)\) lie in \(\mathscr R_b(A)\) by clause~(i)
for the second interpretation, so by (F2) the first guard holds of them and
every value of \(\theta\) is a value of \(\theta_1\) at a true guard, hence in
\(\mathscr R_a(A)\); and since the coding tuples \(\bar w_s\) reach every
member of \(\mathscr R_b(A)\), every member of \(\mathscr R_a(A)\) is a value
of \(\theta\). Guards of both interpretations exist on every infinite set
(first paragraph of the proof of \cref{thm:p03-w03-p03-046-b-transfer}), so
clause~(ii) holds with domain the whole set.

\emph{Involutions.} \(Q_{\mathrm{inv}}\le Q_{1-1}\) with full domain
(\cref{thm:p03-w03-p03-046-b-permutations}) and \(Q_{1-1}\le Q_\phi\) with
full domain by hypothesis, so by composition there is a full-domain
interpretation \((\chi_\iota,x_0=x_0,\iota)\) of \(Q_{\mathrm{inv}}\) in
\(Q_\phi\), with \(\iota(x,y;\bar w_0,\bar w)\): in every infinite set
\(A\) guard tuples \(\bar w_0\) exist, and for each of them the relations
\(\iota(A;\bar w_0,\bar w)\), as \(\bar w\) ranges over points and members
of \(\mathscr R_\phi(A)\), are exactly the graphs of involutions of \(A\).

\emph{Part~(2).} Let \((\chi,\theta_0,\theta)\) interpret \(Q_\psi\) in
\(Q_\phi\), with \(r\) \(k\)-ary. For involution variables
\(b_1,b_2,c\) write \(\mathrm{Mv}_b(x)\) for \(\neg b(x,x)\) and put
\begin{align*}
F(x,y;b_1,b_2,c)&:=\bigl(\mathrm{Mv}_{b_1}(x)\wedge\neg\mathrm{Mv}_c(x)\wedge
b_1(x,y)\bigr)\vee\bigl(\neg(\mathrm{Mv}_{b_1}(x)\wedge\neg\mathrm{Mv}_c(x))
\wedge b_2(x,y)\bigr),\\
\mathrm{Bij}(\bar y_0,\bar s_0;b_1,b_2,c)&:=\forall x\,\exists y\,\bigl(F(x,y)
\wedge\forall y'\,(F(x,y')\to y'=y)\bigr)\wedge\forall x\,\forall x'\,\forall
y\,\bigl(F(x,y)\wedge F(x',y)\to x=x'\bigr)\\
&\qquad\wedge\forall y\,\bigl(\theta_0(y,\bar y_0,\bar s_0)\leftrightarrow
\exists x\,F(x,y)\bigr).
\end{align*}
\(\mathrm{Bij}\) says that \(F\) is the graph of a bijection from the whole
set onto the domain \(\theta_0(\cdot,\bar y_0,\bar s_0)\). Let
\(F^\iota\), \(\mathrm{Bij}^\iota\) arise by replacing each atom
\(b_1(t,t')\), \(b_2(t,t')\), \(c(t,t')\) by
\(\iota(t,t';\bar w_0,\bar w_{b_1})\), \(\iota(t,t';\bar w_0,\bar
w_{b_2})\), \(\iota(t,t';\bar w_0,\bar w_c)\) with fresh tuples. The new
interpretation has guard variables
\(\bar y_0,\bar s_0,\bar w_0,\bar w_{b_1},\bar w_{b_2},\bar w_c\) and
\begin{align*}
\chi^*&:=\chi(\bar y_0,\bar s_0)\wedge\chi_\iota(\bar w_0)\wedge
\mathrm{Bij}^\iota,\qquad\theta_0^*(x_0):=x_0=x_0,\\
\theta^*(x_0,\dots,x_{k-1};\dots,\bar y,\bar s)&:=\exists
z_0\cdots\exists z_{k-1}\Bigl(\textstyle\bigwedge_{i<k}F^\iota(x_i,z_i)\wedge
\theta(z_0,\dots,z_{k-1},\bar y_0,\bar s_0,\bar y,\bar s)\Bigr).
\end{align*}

\emph{Clause~(i).} Suppose \(A\entails\chi^*\) at given parameters. By
the choice of \(\iota\) and (F2) the three relations
\(\beta_1,\beta_2,\gamma\) defined by \(\iota\) at these parameters are
involutions and \(\mathrm{Bij}\) holds of them, so \(F\) defines a bijection
\(f\colon A\to D\), \(D=\theta_0(A,\bar a_0,\bar S_0)\). The domain of the
new interpretation is \(A\), which is infinite because it contains the
infinite set \(D\). By clause~(i) of the given interpretation the values
\(\theta(A,\bar a_0,\bar S_0,\bar a,\bar S)\) are exactly the members of
\(\mathscr R_\psi(D)\), each a relation on \(D\). The value of
\(\theta^*\) at \(\bar a,\bar S\) is
\(\{\bar x\in A^k:(fx_0,\dots,fx_{k-1})\in\theta(A,\bar a_0,\bar S_0,\bar
a,\bar S)\}=f^{-1}[\theta(A,\bar a_0,\bar S_0,\bar a,\bar S)]\). By (F1),
applied to the bijection \(f^{-1}\colon D\to A\), \(R'\mapsto f^{-1}[R']\)
maps \(\mathscr R_\psi(D)\) onto \(\mathscr R_\psi(A)\). So the values of
\(\theta^*\) are exactly the members of \(\mathscr R_\psi(A)\).

\emph{Clause~(ii).} Let \(A\) be infinite, \(\kappa=|A|\). By part~(1)
choose guard parameters \(\bar a_0,\bar S_0\) of the given interpretation
in \(A\) with \(|D|=\kappa\), \(D=\theta_0(A,\bar a_0,\bar S_0)\).
Partition \(D\) into \(D_1,D_2\) with \(|D_1|=|D_2|=\kappa\) (possible as
\(\kappa\cdot2=\kappa\)), and put \(A_1=D_2\cup(A-D)\), \(A_2=D_1\). Then
\(A=A_1\cup A_2\) disjointly, \(A_1\cap D_1=\varnothing\),
\(A_2\cap D_2=\varnothing\), and all four sets have cardinality \(\kappa\).
Choose bijections \(f_1\colon A_1\to D_1\) and \(f_2\colon A_2\to D_2\). Let
\(\beta_1\) exchange \(x\) and \(f_1(x)\) for each \(x\in A_1\) (it moves
every point, since \(A_1\cup D_1=A\) and \(A_1\cap D_1=\varnothing\)), let
\(\beta_2\) exchange \(x\) and \(f_2(x)\) for each \(x\in A_2\) and fix
\(A-D\), and let \(\gamma\) be an involution with
\(\operatorname{Mov}(\gamma)=D_1\), obtained by splitting the infinite set
\(D_1\) into two-element blocks. For \(x\in A_1\): \(x\) is moved by
\(\beta_1\) and not by \(\gamma\), so \(F\) sends \(x\) to
\(\beta_1(x)=f_1(x)\). For \(x\in A_2=D_1\): \(x\) is moved by \(\gamma\),
so \(F\) sends \(x\) to \(\beta_2(x)=f_2(x)\). Thus \(F\) defines
\(f_1\cup f_2\), a bijection from \(A\) onto \(D_1\cup D_2=D\), and
\(\mathrm{Bij}\) holds. A guard tuple \(\bar w_0\) for \(\chi_\iota\) exists
in \(A\), and coding tuples give \(\beta_1,\beta_2,\gamma\) as values of
\(\iota\). By (F2) \(\chi^*\) holds at these parameters, and the domain is
\(A\) itself.
\end{proof}
\end{annexitem}

\begin{annexitem}{annex:p03-w03-p03-046-b-transfer-general}{thm:p03-w03-p03-046-b-transfer-general}{Transfer and Hanf numbers for quantifiers coding pairings}
\begin{proof}
\(Q_{1-1}\le Q_{1-1}\) with full domain by the identity interpretation (no
guard, \(\theta_0:=x_0=x_0\), \(\theta(x,y;s):=s(x,y)\)).
\(Q_{1-1}\le Q_{\mathrm{inv}}\le Q_{\mathrm{perm}}\) with full domain by
\cref{thm:p03-w03-p03-046-b-permutations} and composition (proof of
\cref{thm:p03-w03-p03-046-b-full-domain}), \(Q_{1-1}\le Q_{\mathrm{II}}\)
with full domain by \cref{thm:p03-w03-p03-046-b-maximal}, and under the
hypothesis of \cref{thm:p03-w03-p03-046-b-perm-crit} with full domain by
that lemma. So each listed \(Q_\phi\) satisfies the hypothesis of
\cref{thm:p03-w03-p03-046-b-full-domain}, and part~(2) of that lemma turns
the given interpretation of \(Q_\psi\) into a full-domain one.
\Cref{thm:p03-w03-p03-046-b-transfer} then gives the transfer statements,
and \cref{thm:p03-w03-p03-046-b-hanf}, whose hypothesis is now met, the
inequalities between Hanf and L\"owenheim numbers.
\end{proof}
\end{annexitem}

\begin{annexitem}{annex:p03-w03-p03-046-b-wellorder}{thm:p03-w03-p03-046-b-wellorder}{Well-orders in monadic and first-order logic of order}
\begin{proof}
(1) The monadic condition is valid, so \(Q_{\mathrm{mon}}r\) ranges over
all subsets of \(M\). The formula under \(\neg Q_{\mathrm{mon}}r\) says that
\(r\) is a nonempty set each member of which has a smaller member in the
set, that is, a nonempty set without a least element. So
\((M,<)\entails\mathrm{WO}\) iff every nonempty subset of \(M\) has a least
element, which for a linear order is the definition of a well-order.

(2) The quantifier \(Q_{\mathrm I}\) adds nothing to first-order logic: a
formula \(Q_{\mathrm I}r\,\beta(r)\) is equivalent to \(\exists
z\,\beta[z/r]\), where each atom \(r(t)\) is replaced by \(t=z\), because
the relations quantified over are exactly the singletons. Applying this
replacement from the inside out turns a sentence of
\((\mathrm{Th}(<),Q_{\mathrm I})\) into an equivalent first-order sentence
\(\sigma\) of \(L(<)\). Suppose \(\sigma\) holds in exactly the infinite
well-orders among the infinite linear orders. Add constants
\(c_0,c_1,\dots\) and let \(\Gamma\) consist of \(\sigma\), the axioms of
linear order, the sentences \(\exists^{\ge n}x\,(x=x)\) for \(n<\omega\),
and \(c_{n+1}<c_n\) for \(n<\omega\). A finite part of \(\Gamma\) mentions
only \(c_0,\dots,c_n\) for some \(n\); it holds in \((\omega,<)\) with
\(c_i\) interpreted as \(n-i\), since \((\omega,<)\) is an infinite
well-order and so satisfies \(\sigma\). By compactness
(\cref{thm:p03-compactness-ultraproduct}) \(\Gamma\) has a model. It is an
infinite linear order satisfying \(\sigma\) in which \(\{c_n:n<\omega\}\)
has no least element, so it is not a well-order: a contradiction.
\end{proof}
\end{annexitem}

\begin{annexitem}{annex:p03-w03-p03-046-b-maximal}{thm:p03-w03-p03-046-b-maximal}{\(Q_{\mathrm{II}}\) is maximal (Proposition~1.2.7)}
\begin{proof}
Let \(r\) be \(k\)-ary, \(k\ge1\). The parameters of \(Q_{\mathrm{II}}\)
are arbitrary binary relations.

\emph{Tuple coding.} Let \(\mathrm{Pair}(p_0,p_1)\) say that \(p_0\) and
\(p_1\) are graphs of total functions and that for all \(x,y\) there is
exactly one \(z\) with \(p_0(z,x)\wedge p_1(z,y)\). Under
\(\mathrm{Pair}(P_0,P_1)\) the map \(z\mapsto(P_0(z),P_1(z))\) is a
bijection \(A\to A\times A\). Put \(\mathrm{Tup}_1(z,x_1):=z=x_1\) and
\[
\mathrm{Tup}_{j+1}(z,x_1,\dots,x_{j+1}):=\exists
w\,\bigl(p_0(z,x_1)\wedge p_1(z,w)\wedge\mathrm{Tup}_j(w,x_2,\dots,x_{j+1})\bigr).
\]
By induction on \(j\), under \(\mathrm{Pair}\) the relation
\(\mathrm{Tup}_j\) is the graph of a bijection \(t_j\colon A^j\to A\): for
\(j=1\) it is the identity, and \(\mathrm{Tup}_{j+1}(z,\bar x)\) holds iff
\(P_0(z)=x_1\) and \(t_j(x_2,\dots,x_{j+1})=P_1(z)\), which is the
composite of the bijections \(A\times A^j\to A\times A\to A\). Let
\[
\mathrm{Dec}(\bar x;p_0,p_1,d):=\exists z\,(\mathrm{Tup}_k(z,\bar x)\wedge
d(z,z)).
\]
Under \(\mathrm{Pair}(P_0,P_1)\), \(\mathrm{Dec}(A;P_0,P_1,D)=\{\bar
x:(t_k(\bar x),t_k(\bar x))\in D\}\), and every \(R\subseteq A^k\) arises,
from \(D=\{(t_k(\bar x),t_k(\bar x)):\bar x\in R\}\). Let
\(\psi^{\mathrm{Dec}}(p_0,p_1,d)\) be \(\psi\) with each atom \(r(\bar u)\)
replaced by \(\mathrm{Dec}(\bar u;p_0,p_1,d)\); by (F2) it holds iff
\(\mathrm{Dec}(A;P_0,P_1,D)\in\mathscr R_\psi(A)\).

\emph{The interpretation.} Guard
\(\chi(p_0,p_1,d_0):=\mathrm{Pair}(p_0,p_1)\wedge\psi^{\mathrm{Dec}}(p_0,p_1,d_0)\);
\(\theta_0:=x_0=x_0\); and, with the coding variable \(d\),
\[
\theta(\bar x;p_0,p_1,d_0;d):=\bigl(\psi^{\mathrm{Dec}}(p_0,p_1,d)\wedge
\mathrm{Dec}(\bar x;p_0,p_1,d)\bigr)\vee\bigl(\neg\psi^{\mathrm{Dec}}(p_0,p_1,d)
\wedge\mathrm{Dec}(\bar x;p_0,p_1,d_0)\bigr).
\]
Clause~(i): if the guard holds in \(A\), then \(\theta_0(A)=A\) is
infinite; every value of \(\theta\) lies in \(\mathscr R_\psi(A)\), in the
first case by \(\psi^{\mathrm{Dec}}(D)\), in the second by the guard; and
every \(R\in\mathscr R_\psi(A)\) is \(\mathrm{Dec}(A;P_0,P_1,D)\) for some
\(D\), for which the first case applies. Clause~(ii): let \(B\) be
infinite and \(A=B\). Since \(|B\times B|=|B|\) there is a bijection
\(B\to B\times B\), whose coordinate functions give \(P_0,P_1\) with
\(\mathrm{Pair}\). By hypothesis the first-order sentence \(\psi(R)\), in
the vocabulary \(\{R\}\), has an infinite model. Adding \(|B|\) new
constants required to be distinct keeps it finitely satisfiable in that
model, so by \cref{thm:p03-compactness-ultraproduct} it has a model with at
least \(|B|\) elements, and by \cref{thm:p03-w03-p03-039-downward-ls}
(for the fragment of first-order formulas) an elementary substructure
containing those constants has exactly \(|B|\) elements; by (F1) there is
then a model with universe \(B\), say \((B,R_0)\). With
\(D_0\) coding \(R_0\) the guard holds, and \(\theta_0(B)=B\).
\end{proof}
\end{annexitem}

\begin{annexitem}{annex:p03-w03-p03-046-b-mon}{thm:p03-w03-p03-046-b-mon}{An infinite, coinfinite set gives monadic logic (Lemma~1.2.8)}
\begin{proof}
Let \(A_1\) be infinite, \(\bar a_1\in A_1\), \(\bar R_1\in\mathscr
R_\psi(A_1)\) (each entry) and \(\phi(x,\bar z,\bar r)\) first-order with
\(\phi(A_1,\bar a_1,\bar R_1)\) and its complement both infinite. Call
\(X\subseteq A\) \emph{regular} if \(|X|=|A-X|=|A|\); the complement of a
regular set is regular.

\emph{Step 1: a definable regular set in each cardinality.} Fix
\(\kappa\ge\aleph_0\). In the vocabulary with relation symbols \(\bar R\),
constants \(\bar c\) and new constants \(d_\alpha,e_\alpha\)
(\(\alpha<\kappa\)), let \(\Gamma\) contain \(\psi(R)\) for each \(R\) in
\(\bar R\), \(\phi(d_\alpha,\bar c,\bar R)\), \(\neg\phi(e_\alpha,\bar
c,\bar R)\), and \(d_\alpha\neq d_\beta\), \(e_\alpha\neq e_\beta\) for
\(\alpha\neq\beta\). A finite part of \(\Gamma\) is satisfied in
\((A_1,\bar a_1,\bar R_1)\) by sending the finitely many \(d\)'s and \(e\)'s
to distinct points of the two infinite sets. By
\cref{thm:p03-compactness-ultraproduct} \(\Gamma\) has a model, and by
\cref{thm:p03-w03-p03-039-downward-ls} (for the fragment of first-order formulas) an elementary substructure of it containing the
\(2\kappa\) constants has cardinality at most \(\kappa\), hence exactly
\(\kappa\). Call it \((A_\kappa,\bar a_\kappa,\bar R_\kappa)\). Then
\(\bar R_\kappa\in\mathscr R_\psi(A_\kappa)\) and
\(\phi(A_\kappa,\bar a_\kappa,\bar R_\kappa)\) is regular.

\emph{Step 2: every regular set.} Let \(|A|=\kappa\) and \(X\subseteq A\)
regular. As \(\phi(A_\kappa,\dots)\), its complement, \(X\) and \(A-X\)
all have cardinality \(\kappa\), there is a bijection \(g\colon A_\kappa\to
A\) with \(g[\phi(A_\kappa,\bar a_\kappa,\bar R_\kappa)]=X\). By (F1),
\(g[\bar R_\kappa]\in\mathscr R_\psi(A)\) and
\(\phi(A,g\bar a_\kappa,g[\bar R_\kappa])=X\).

\emph{Step 3: Boolean combinations.} Let \(X\subseteq A\), \(|A|=\kappa\).
(a) If \(|A-X|=\kappa\) and \(|X|<\kappa\), split \(A-X=W\cup W'\) into
disjoint sets of cardinality \(\kappa\); then \(Y=X\cup W\) and \(A-W\) are
regular and \(X=Y\cap(A-W)\). (b) If \(|A-X|=|X|=\kappa\), split
\(A-X=W\cup W'\) as before; \(X\) and \(A-W\supseteq X\cup W'\) are
regular and \(X=X\cap(A-W)\). (c) If \(|A-X|<\kappa\), then
\(|X|=\kappa\), and by (a) applied to \(S=A-X\) there are regular
\(Y,Y'\) with \(S=Y\cap Y'\); so \(X=(A-Y)\cup(A-Y')\), a union of two
regular sets. In every case \(X=(Y_1\cap Y_2)\cup(Y_3\cap Y_4)\) with
\(Y_1,\dots,Y_4\) regular (repeat a set where needed). Put
\[
\mu(x;\bar z_1,\bar r_1,\dots,\bar z_4,\bar r_4):=
(\phi(x,\bar z_1,\bar r_1)\wedge\phi(x,\bar z_2,\bar r_2))\vee
(\phi(x,\bar z_3,\bar r_3)\wedge\phi(x,\bar z_4,\bar r_4)).
\]
By Step~2 the sets \(\mu(A,\dots)\) with parameters from \(A\) and
\(\mathscr R_\psi(A)\) are exactly the subsets of \(A\).

\emph{Step 4: the interpretation.} No guard parameters are needed: \(\chi\)
is the true formula, \(\theta_0:=x_0=x_0\) and \(\theta:=\mu\). The
relations of \(Q_{\mathrm{mon}}\) on \(A\) are all subsets of \(A\), so
clause~(i) is Step~3, and clause~(ii) holds with \(A=B\).

\emph{Converse.} Suppose \(Q_{\mathrm{mon}}\le Q_\psi\) by
\((\chi,\theta_0,\theta)\). Clause~(ii) with \(B=\omega\) gives \(A\),
\(\bar a_0\), \(\bar S_0\) with the guard true; \(D=\theta_0(A,\bar
a_0,\bar S_0)\) is infinite. Choose \(X\subseteq D\) with \(X\) and
\(D-X\) infinite. By clause~(i), \(X=\theta(A,\bar a_0,\bar S_0,\bar
a,\bar S)\) for some \(\bar a\in A\) and \(\bar S\in\mathscr R_\psi(A)\);
as \(\bar S_0\in\mathscr R_\psi(A)\) too, \(X\) is simply definable by
\(Q_\psi\), and it is infinite with \(A-X\supseteq D-X\) infinite.
\end{proof}
\end{annexitem}

\begin{annexitem}{annex:p03-w03-p03-046-b-uniform}{thm:p03-w03-p03-046-b-uniform}{Strong minimality and uniform bounds}
\begin{proof}
\emph{Strong minimality.} If some \(\phi(A,\bar b,\bar R)\) with \(\bar
R\in\mathscr R_\psi(A)\) were infinite and coinfinite, \(Q_\psi\) would
simply define an infinite, coinfinite set, and
\cref{thm:p03-w03-p03-046-b-mon} would give \(Q_{\mathrm{mon}}\le Q_\psi\).

\emph{Uniform bound.} Suppose every \(\phi(A,\bar b,\bar R)\) is finite
but their sizes are unbounded. Let \(\Gamma\) be the elementary diagram of
\((A,\bar R)\) (a constant for each element of \(A\)) together with
\(\exists^{\ge n}x\,\phi(x,\bar c,\bar R)\) for all \(n\), where \(\bar
c\) are new constants and \(\exists^{\ge n}\) is the negation of
\(\exists^{<n-1}\). A finite part of \(\Gamma\) mentions some largest
\(n\); interpreting \(\bar c\) by a \(\bar b\) with
\(|\phi(A,\bar b,\bar R)|\ge n\) satisfies it in \((A,\bar R)\). By
\cref{thm:p03-compactness-ultraproduct} \(\Gamma\) has a model
\((A^*,\bar R^*,\bar c^*)\); since it satisfies the elementary diagram,
\(a\mapsto a^{A^*}\) is an elementary embedding of \((A,\bar R)\) into
\((A^*,\bar R^*)\). So \(A^*\) is infinite and \(\psi(\bar R^*)\) holds,
that is \(\bar R^*\in\mathscr R_\psi(A^*)\), and
\(\phi(A^*,\bar c^*,\bar R^*)\) is infinite. By the first part applied to
\((A^*,\bar R^*)\) it is cofinite; let \(m=|A^*-\phi(A^*,\bar c^*,\bar
R^*)|\). Then \((A^*,\bar R^*)\entails\exists\bar
y\,(\exists^{<m}x)\neg\phi(x,\bar y,\bar R)\), so the same holds in
\((A,\bar R)\): some \(\phi(A,\bar b,\bar R)\) has a complement of at most
\(m\) points and is therefore infinite, a contradiction.
\end{proof}
\end{annexitem}

\begin{annexitem}{annex:p03-w03-p03-046-b-perm-crit}{thm:p03-w03-p03-046-b-perm-crit}{An equivalence relation gives permutational logic (Lemma~1.2.9)}
\begin{proof}
Let \(\varepsilon(x,y,\bar a_1,\bar R_1)\), \(\bar R_1\in\mathscr
R_\psi(A_1)\), define on the infinite set \(A_1\) an equivalence relation
\(E\) with infinitely many infinite classes. Since
\(Q_{\mathrm{mon}}\le Q_\psi\), the converse part and Step~3 of
\cref{thm:p03-w03-p03-046-b-mon} give a formula \(\mu(x,\bar w)\) (the
tuple \(\bar w\) containing individual and relation variables) whose
values on any infinite set, with relation parameters from \(\mathscr
R_\psi\), are exactly its subsets.

\emph{Step 1: a definable matching.} Choose distinct classes \(K_i\)
(\(i<\omega\)) and \(u_i\neq v_i\) in \(K_i\), and put
\(B_1=\{u_i,v_i:i<\omega\}\). Then \(A_1-B_1\supseteq K_0-\{u_0,v_0\}\) is
infinite. Fix parameters \(\bar p_1\) with \(\mu(A_1,\bar p_1)=B_1\) and put
\[
\pi(x,y;\bar z,\bar r,\bar w):=\mu(x,\bar w)\wedge\mu(y,\bar w)\wedge
x\neq y\wedge\varepsilon(x,y,\bar z,\bar r).
\]
For distinct \(x,y\in B_1\), \(\varepsilon\) holds iff \(\{x,y\}=\{u_i,v_i\}\)
for some \(i\); so \(\pi(A_1,\bar a_1,\bar R_1,\bar p_1)\) pairs each
\(u_i\) with \(v_i\). Let \(\mathrm{Mt}(\bar z,\bar r,\bar w)\) say that
\(\pi\) is symmetric and that every \(x\) with \(\mu(x,\bar w)\) has
exactly one \(y\) with \(\pi(x,y)\); since \(\pi\) is irreflexive and only
relates points of \(\mu\), \(\mathrm{Mt}\) says that \(\pi\) is a perfect
matching of the set defined by \(\mu\).

\emph{Step 2: every cardinality.} Fix \(\kappa\ge\aleph_0\). As in Step~1 of
\cref{thm:p03-w03-p03-046-b-mon}, compactness applied to \(\psi\) of each
relation parameter, \(\mathrm{Mt}\), \(\kappa\) distinct new constants
satisfying \(\mu\) and \(\kappa\) satisfying \(\neg\mu\) (finitely
satisfiable in \(A_1\), where \(B_1\) and \(A_1-B_1\) are infinite), followed
by \cref{thm:p03-w03-p03-039-downward-ls} (for the fragment of first-order formulas), gives a set \(A_\kappa\) of cardinality
\(\kappa\) with parameters \(\bar q_\kappa\) (relations in \(\mathscr
R_\psi(A_\kappa)\)) such that \(\mu(A_\kappa,\bar q_\kappa)\) is regular
and \(\pi(A_\kappa,\bar q_\kappa)\) is a perfect matching of it.

\emph{Step 3: a definable pairing of every infinite set.} Let \(|A|=\kappa\)
and \(A=X\cup X'\) with \(X,X'\) disjoint of cardinality \(\kappa\). As in
Step~2 of \cref{thm:p03-w03-p03-046-b-mon}, bijections \(g,g'\colon
A_\kappa\to A\) carrying \(\mu(A_\kappa,\bar q_\kappa)\) onto \(X\),
respectively \(X'\), give parameters \(\bar q=g\bar q_\kappa\),
\(\bar q'=g'\bar q_\kappa\) for which \(\pi(A,\bar q)\) is a perfect
matching of \(X\) and \(\pi(A,\bar q')\) one of \(X'\), by (F1). Then
\[
\rho(x,y;\bar q,\bar q'):=x=y\vee\pi(x,y;\bar q)\vee\pi(x,y;\bar q')
\]
defines an equivalence relation on \(A\) every class of which has exactly
two elements.

\emph{Step 4: every pairing.} Any two such equivalence relations \(P,P'\)
on \(A\) have \(|A|\) classes each, so a bijection between their sets of
classes, combined with bijections between corresponding classes, gives a
permutation \(h\) of \(A\) with \(h[P]=P'\); by (F1)
\(\rho(A,h\bar q,h\bar q')=P'\).

\emph{The interpretation.} Let \(\mathrm{Pg}(\bar q,\bar q')\) say that
\(\rho\) is an equivalence relation all of whose classes have exactly two
elements. Guard \(\chi(\bar q_0,\bar q_0'):=\mathrm{Pg}(\bar q_0,\bar
q_0')\); \(\theta_0:=x_0=x_0\);
\[
\theta(x,y;\bar q_0,\bar q_0';\bar q,\bar q'):=
(\mathrm{Pg}(\bar q,\bar q')\wedge\rho(x,y;\bar q,\bar q'))\vee
(\neg\mathrm{Pg}(\bar q,\bar q')\wedge\rho(x,y;\bar q_0,\bar q_0')).
\]
Every value is a pairing, and by Steps~3 and~4 every pairing of \(A\) is a
value; this is clause~(i). Clause~(ii) holds with \(A=B\), the guard
parameters coming from Step~3.
\end{proof}
\end{annexitem}

\begin{annexitem}{annex:p03-w03-p03-046-b-permutations}{thm:p03-w03-p03-046-b-permutations}{Permutations, involutions and pairings}
\begin{proof}
\emph{Factorization.} Let \(f\) be a permutation of \(A\). Choose a point
\(x_0\) in each orbit and index the orbit by \(x_i=f^i(x_0)\), with
\(i\in\mathbb Z\) if the orbit is infinite and \(i\in\mathbb Z/n\mathbb Z\)
if it has \(n\) elements. On the orbit put \(\tau(x_i)=x_{-i}\) and
\(\sigma(x_i)=x_{1-i}\). These are well defined (also modulo \(n\)) and
are involutions, since \(-(-i)=i\) and \(1-(1-i)=i\), and
\(\sigma(\tau(x_i))=x_{1+i}=f(x_i)\). Taking the union over the orbits,
\(f=\sigma\circ\tau\).

\(Q_{\mathrm{perm}}\le Q_{\mathrm{inv}}\): no guard, \(\theta_0:=x_0=x_0\),
\(\theta(x,y;\sigma,\tau):=\exists z\,(\tau(x,z)\wedge\sigma(z,y))\). A
composite of two permutations is a permutation, and every permutation is
such a composite.

\(Q_{\mathrm{inv}}\le Q_{\mathrm{perm}}\): no guard,
\(\theta(x,y;f):=(\mathrm{Sym}(f)\wedge f(x,y))\vee(\neg\mathrm{Sym}(f)\wedge
x=y)\), where \(\mathrm{Sym}(f):=\forall x\forall y\,(f(x,y)\to f(y,x))\).
A permutation equals its inverse iff its graph is symmetric, and the
identity is an involution.

\(Q_{1-1}\le Q_{\mathrm{inv}}\): let \(\mathrm{Fpf}(\alpha):=\forall
x\,\neg\alpha(x,x)\); guard \(\chi(\alpha_0):=\mathrm{Fpf}(\alpha_0)\), and
\(\theta(x,y;\alpha_0;\alpha):=(\mathrm{Fpf}(\alpha)\wedge(x=y\vee\alpha(x,y)))\vee
(\neg\mathrm{Fpf}(\alpha)\wedge(x=y\vee\alpha_0(x,y)))\). For a
fixed-point-free involution \(\alpha\) the relation \(x=y\vee\alpha(x,y)\)
is a pairing, and every pairing \(P\) arises from the fixed-point-free
involution exchanging the two points of each class. An infinite set has a
fixed-point-free involution, since \(|B\times2|=|B|\); this gives
clause~(ii).

\(Q_{\mathrm{inv}}\le Q_{1-1}\): \(Q_{1-1}\) simply defines an infinite,
coinfinite set (\cref{def:p03-w03-p03-046-b-example} gives one on
\(\omega\)), so \cref{thm:p03-w03-p03-046-b-mon} gives \(\mu\) defining
every subset from pairings and points. With \(\mu_F,\mu_D,\mu_{D'}\)
three copies of \(\mu\) with separate parameters and \(E,E'\) pairing
variables, put
\[
\iota(x,y):=(\mu_F(x)\wedge y=x)\vee(\mu_D(x)\wedge E(x,y)\wedge x\neq y)
\vee(\mu_{D'}(x)\wedge E'(x,y)\wedge x\neq y),
\]
let \(\mathrm{Inv}(\iota)\) say that \(\iota\) is the graph of a total
function equal to its own inverse, and let
\(\theta:=(\mathrm{Inv}(\iota)\wedge\iota)\vee(\neg\mathrm{Inv}(\iota)\wedge
x=y)\); there is no guard. Every value is an involution. Let \(\alpha\) be
an involution of the infinite set \(A\), with fixed set \(F\) and moved
set \(N\); \(N\) is a union of two-element \(\alpha\)-orbits. If \(N\) is
infinite, split its orbits into two infinite families with unions \(D\),
\(D'\); if \(N\) is finite, let \(D=N\) and \(D'=\varnothing\), so that
\(F\) is infinite. In both cases \(A-D\) and \(A-D'\) are infinite. Let
\(E\) agree with \(\alpha\) on \(D\) and pair \(A-D\) arbitrarily, and let
\(E'\) agree with \(\alpha\) on \(D'\) and pair \(A-D'\) arbitrarily; an
infinite set has a pairing. Since \(F,D,D'\) partition \(A\), \(\iota\)
with these parameters is the graph of \(\alpha\), so \(\mathrm{Inv}\)
holds and \(\theta\) gives \(\alpha\). This is clause~(i); clause~(ii)
holds with \(A=B\).
\end{proof}
\end{annexitem}

\begin{annexitem}{annex:p03-w03-p03-046-b-mso-sets}{thm:p03-w03-p03-046-b-mso-sets}{Monadic logic of unary structures}
\begin{proof}
Let \(L=\{P_1,\dots,P_m\}\). Formulas have individual variables and set
variables, and the quantifier rank counts both kinds of quantifier. Put
\(t_0=0\) and \(t_{n+1}=2t_n+1\). For a structure \(\mathfrak A=(A,\bar
P)\), points \(\bar a=(a_1,\dots,a_p)\) and sets \(\bar
S=(S_1,\dots,S_q)\), and \(\varepsilon\in\{0,1\}^{m+q}\), let
\(A_\varepsilon\) be the set of points whose membership pattern in
\(P_1,\dots,P_m,S_1,\dots,S_q\) is \(\varepsilon\), and
\(A^-_\varepsilon=A_\varepsilon-\{a_1,\dots,a_p\}\). Call
\((\mathfrak A,\bar a,\bar S)\) and \((\mathfrak B,\bar b,\bar T)\)
\emph{\(n\)-similar} if (1) \(a_i=a_j\iff b_i=b_j\); (2) \(a_i\) and
\(b_i\) have the same pattern; (3)
\(\min(|A^-_\varepsilon|,t_n)=\min(|B^-_\varepsilon|,t_n)\) for every
\(\varepsilon\).

\emph{Claim.} \(n\)-similar triples satisfy the same formulas of rank at
most \(n\). By induction on the formula; \(n\)-similarity implies
\(n'\)-similarity for \(n'\le n\) because \(t_{n'}\le t_n\). Atoms
\(x_i=x_j\), \(P(x_i)\), \(X_j(x_i)\) are settled by (1) and (2), and
connectives are immediate. Since similarity is symmetric and \(\forall\)
is \(\neg\exists\neg\), it suffices to treat \(\exists x\,\sigma\) and
\(\exists X\,\sigma\) of rank at most \(n\), \(n\ge1\), true on the left,
and to find a witness on the right making the extended triples
\((n-1)\)-similar.

\emph{Point.} Let \(a\) witness. If \(a=a_i\), take \(b=b_i\); the cells are
unchanged. Otherwise \(a\in A^-_\varepsilon\) for some \(\varepsilon\), so
\(|A^-_\varepsilon|\ge1\) and, as \(t_n\ge1\),
\(B^-_\varepsilon\neq\varnothing\); take \(b\in B^-_\varepsilon\). Then
(1), (2) hold, and the cell \(\varepsilon\) loses one point on each side:
if both sizes were below \(t_n\) they stay equal, and if both were at
least \(t_n\) both become at least \(t_n-1=2t_{n-1}\ge t_{n-1}\).

\emph{Set.} Let \(S\) witness. Each cell \(A^-_\varepsilon\) splits into
\(A^-_\varepsilon\cap S\) and \(A^-_\varepsilon-S\), of sizes \(s_1,s_2\).
Split \(B^-_\varepsilon\) into two parts of sizes \(s_1',s_2'\) with
\(\min(s_i',t_{n-1})=\min(s_i,t_{n-1})\): if \(|A^-_\varepsilon|<t_n\) the two
cells have equal finite size and the split is copied; if both have at least
\(t_n\) points and \(s_1<t_{n-1}\), take \(s_1'=s_1\), and then
\(s_2,s_2'\ge t_n-t_{n-1}>t_{n-1}\); symmetrically if \(s_2<t_{n-1}\); if
\(s_1,s_2\ge t_{n-1}\), give the first part \(t_{n-1}\) points, leaving at
least \(t_n-t_{n-1}\ge t_{n-1}\) for the second. Let \(T\) be the union of
the first parts together with \(\{b_i:a_i\in S\}\). Then (1)--(3) hold at
rank \(n-1\) for the extended triples.

\emph{Part~1.} Let \(\sigma\) be a sentence of rank \(n\) with an infinite
model \(\mathfrak A\). Some cell \(A_\varepsilon\) (no parameters) is
infinite, since there are finitely many. For \(\lambda\ge\aleph_0\) let
\(\mathfrak B\) have cells \(B_\varepsilon\) with
\(|B_\varepsilon|=|A_\varepsilon|\) if \(A_\varepsilon\) is finite and
\(|B_\varepsilon|=\lambda\) otherwise (a disjoint union, with \(P_j\)
the union of the cells whose pattern has \(1\) at \(j\)). Then
\(\lvert\dom\mathfrak B\rvert=\lambda\), the structures are \(n\)-similar, and
\(\mathfrak B\entails\sigma\). With \(\lambda=\aleph_0\) this bounds the
L\"owenheim number by \(\aleph_0\); with arbitrary \(\lambda\) every
sentence with an infinite model has arbitrarily large models, so the Hanf
number is \(\aleph_0\).

\emph{Part~2.} With \(m=0\) and no parameters there is one cell, the
whole set. Two infinite sets are therefore \(n\)-similar for every \(n\),
so they satisfy the same sentences of monadic logic, which are the
sentences of \((\mathrm{Th}(=),Q_{\mathrm{mon}})\).
\end{proof}
\end{annexitem}

\begin{annexitem}{annex:p03-w03-p03-046-b-fv}{thm:p03-w03-p03-046-b-fv}{A Feferman--Vaught theorem for unary monadic logic}
\begin{proof}
Write \(t=t_n\) and \(\tau(x)=\min(x,t)\) for cardinals \(x\). For an
\(L\)-structure \(\mathfrak A\) and \(\varepsilon\in\{0,1\}^m\), \(A_\varepsilon\)
is the set of points whose membership pattern in \(P_1,\dots,P_m\) is
\(\varepsilon\), and the \emph{profile} of \(\mathfrak A\) is
\(c(\mathfrak A)=(\tau(|A_\varepsilon|))_{\varepsilon}\), a member of the
finite set \(V=\{0,\dots,t\}^{\{0,1\}^m}\). For \(c\in V\) let \(\delta_c\) be
the conjunction over \(\varepsilon\) of ``exactly \(c_\varepsilon\) points
have pattern \(\varepsilon\)'' if \(c_\varepsilon<t\) and ``at least \(t\)
points have pattern \(\varepsilon\)'' if \(c_\varepsilon=t\). Written with
\(\exists^{\ge j}\) as a block of \(j\) existential quantifiers, ``at least
\(t\)'' has rank \(t\), and ``exactly \(j\)'' is \(\exists^{\ge
j}\wedge\neg\exists^{\ge j+1}\), of rank \(j+1\le t\) as \(j<t\). So each
\(\delta_c\) is a first-order sentence of rank at most \(t\), and
\(\mathfrak A\entails\delta_c\) iff \(c(\mathfrak A)=c\). Enumerate the
\(\delta_c\) as \(\delta_1,\dots,\delta_N\).

\emph{Step 1: the profile decides rank-\(n\) sentences.} If
\(c(\mathfrak A)=c(\mathfrak A')\), then \(\mathfrak A\) and \(\mathfrak
A'\), with no points and no sets as parameters, are \(n\)-similar in the
sense of the proof of \cref{thm:p03-w03-p03-046-b-mso-sets}: conditions
(1) and (2) are vacuous and (3) is the equality of profiles. By the Claim
in that proof they satisfy the same sentences of rank at most \(n\).

\emph{Step 2: disjoint unions.} A point of \(\mathfrak A\sqcup\mathfrak B\)
has the same pattern as in its summand, so
\((A\sqcup B)_\varepsilon=A_\varepsilon\sqcup B_\varepsilon\). For cardinals
\(x,y\), \(\tau(x+y)=\tau(\tau(x)+\tau(y))\): if \(x,y<t\) the arguments
agree, and if, say, \(x\ge t\), both sides equal \(t\). So
\(c(\mathfrak A\sqcup\mathfrak B)\) is the function \(c\oplus d\),
\((c\oplus d)_\varepsilon=\tau(c_\varepsilon+d_\varepsilon)\), of
\(c=c(\mathfrak A)\), \(d=c(\mathfrak B)\).

\emph{Step 3: products.} The pattern of \((a,b)\) in
\(\mathfrak A\times\mathfrak B\) is the coordinatewise minimum
\(\varepsilon\wedge\varepsilon'\) of the patterns \(\varepsilon\) of \(a\)
and \(\varepsilon'\) of \(b\), so \(|(A\times B)_\eta|=\sum\{|A_\varepsilon|\cdot
|B_{\varepsilon'}|:\varepsilon\wedge\varepsilon'=\eta\}\). For finitely many
pairs of cardinals \((x_i,y_i)\),
\(\tau(\sum_ix_iy_i)=\tau(\sum_i\tau(x_i)\tau(y_i))\). Indeed
\(\tau(x_i)\tau(y_i)\le x_iy_i\) for each \(i\). If equality holds for
every \(i\), the two sides agree. Otherwise for some \(i\) one of
\(x_i,y_i\) is at least \(t\) and neither is \(0\) (a zero factor makes both
products \(0\)); then \(\tau(x_i)\tau(y_i)\ge t\cdot1\), so both sums are at
least \(t\) and both sides equal \(t\). So \(c(\mathfrak A\times\mathfrak B)\)
is a function \(c\otimes d\) of the two profiles.

\emph{Step 4: the reduction.} Let \(\Gamma^\sqcup_\sigma\) be the set of
pairs \((c,d)\in V^2\) for which some structure of profile \(c\oplus d\)
satisfies \(\sigma\), and \(\Gamma^\times_\sigma\) the same with
\(c\otimes d\). If \(\mathfrak A\sqcup\mathfrak B\entails\sigma\), then
\((c(\mathfrak A),c(\mathfrak B))\in\Gamma^\sqcup_\sigma\) by Step~2.
Conversely, if \((c(\mathfrak A),c(\mathfrak B))\in\Gamma^\sqcup_\sigma\),
some structure with the profile \(c(\mathfrak A)\oplus c(\mathfrak
B)=c(\mathfrak A\sqcup\mathfrak B)\) satisfies \(\sigma\), and by Step~1 so
does \(\mathfrak A\sqcup\mathfrak B\). The product case is the same with
Step~3. The final assertion follows because \(\mathfrak A\) and
\(\mathfrak A'\) satisfy the same \(\delta_i\) when they satisfy the same
monadic sentences, every first-order sentence being one.
\end{proof}
\end{annexitem}

\begin{annexitem}{annex:p03-w03-p03-046-b-uncountable}{thm:p03-w03-p03-046-b-uncountable}{A permutational sentence with only uncountable models}
\begin{proof}
Let \(\Sigma\) be the sentence of \(L(Q_{\mathrm{perm}},Q_{\mathrm{mon}})\)
in pure identity theory
\[
Q_{\mathrm{perm}}f\;\neg Q_{\mathrm{mon}}X\,\Bigl(\exists x\,X(x)\wedge
\forall x\bigl(X(x)\leftrightarrow\exists y\,(f(x,y)\wedge X(y))\bigr)\wedge
\exists x\,\neg X(x)\Bigr),
\]
which says that some permutation \(f\) has no invariant subset other than
\(\varnothing\) and the whole set. If an infinite set \(A\) satisfies it
with \(f\), then for any \(a\) the orbit \(\{f^i(a):i\in\mathbb Z\}\) is
nonempty and invariant (\(x\) lies in it iff \(f(x)\) does), so it is all
of \(A\), and \(A\) is countable. If \(A\) is countably infinite, identify
it with \(\mathbb Z\) and let \(f(i)=i+1\); a nonempty invariant set
contains some \(i\) and then every \(i+k\), \(k\in\mathbb Z\), so it is
\(A\). Thus the infinite models of \(\Sigma\) are the countable sets.

By \cref{thm:p03-w03-p03-046-b-permutations}, \(Q_{\mathrm{perm}}\le
Q_{\mathrm{inv}}\le Q_{1-1}\) by full-domain interpretations, and by
\cref{thm:p03-w03-p03-046-b-mon} (with
\cref{def:p03-w03-p03-046-b-example}) so is \(Q_{\mathrm{mon}}\le
Q_{1-1}\). Three applications of
\cref{thm:p03-w03-p03-046-b-transfer}, removing \(Q_{\mathrm{perm}}\) with
\(\mathcal Q=\{Q_{\mathrm{mon}},Q_{\mathrm{inv}}\}\), then
\(Q_{\mathrm{inv}}\) with \(\mathcal Q=\{Q_{\mathrm{mon}},Q_{1-1}\}\), then
\(Q_{\mathrm{mon}}\) with \(\mathcal Q=\{Q_{1-1}\}\), give a sentence
\(\Sigma'\) of \((\mathrm{Th}(=),Q_{1-1})\) equivalent to \(\Sigma\) on every
infinite set. The infinite models of \(\neg\Sigma'\) are the uncountable
sets.
\end{proof}
\end{annexitem}

\begin{annexitem}{annex:p03-w03-p03-046-b-distinct}{thm:p03-w03-p03-046-b-distinct}{The chain and its first two separations}
\begin{proof}
\emph{The chain.} \(Q_{\mathrm I}\le Q_{\mathrm{mon}}\): no guard,
\(\theta_0:=x_0=x_0\), and \(\theta(x;y):=x=y\) with an individual coding
parameter \(y\); the values are exactly the singletons. \(Q_{\mathrm{mon}}\le
Q_{1-1}\) is \cref{thm:p03-w03-p03-046-b-mon} with
\cref{def:p03-w03-p03-046-b-example}. \(Q_{1-1}\le Q_{\mathrm{II}}\) is
\cref{thm:p03-w03-p03-046-b-maximal}, since \(\omega\) has a pairing.

\emph{\(Q_{\mathrm{mon}}\not\le Q_{\mathrm I}\).} Suppose
\((\chi,\theta_0,\theta)\) interprets \(Q_{\mathrm{mon}}\) in
\(Q_{\mathrm I}\); its relation parameters are singletons. Clause~(ii) with
\(B=\omega\) gives \(A,\bar a_0,\bar S_0\) with the guard true. Regard
\(\mathfrak A=(A,\bar a_0,\bar S_0)\) as a first-order structure with
constants for \(\bar a_0\) and unary predicates for \(\bar S_0\). By
\cref{thm:p03-w03-p03-039-downward-ls} (for the fragment of first-order formulas) there is a countably infinite
\(\mathfrak A'\preccurlyeq\mathfrak A\) containing \(\bar a_0\) and the
members of the singletons \(\bar S_0\). Its predicates are the same
singletons, and \(\mathfrak A'\entails\chi(\bar a_0,\bar S_0)\) by
elementarity. Clause~(i) in the countable set \(A'\) makes
\(D=\theta_0(A',\bar a_0,\bar S_0)\) infinite and every subset of \(D\) a
set \(\theta(A',\bar a_0,\bar S_0,\bar a,\bar S)\) with \(\bar a\) a finite
tuple of points and \(\bar S\) a finite tuple of singletons of \(A'\).
There are countably many such parameter tuples, but \(D\) has \(2^{\aleph_0}\)
subsets, a contradiction.

\emph{\(Q_{\mathrm{mon}}\not\le Q_{\mathrm I}\), along Baldwin's route.}
Suppose \(Q_{\mathrm{mon}}\le Q_{\mathrm I}\). By
\cref{thm:p03-w03-p03-046-b-qi-normal}(2) there is a full-domain
interpretation, so by \cref{thm:p03-w03-p03-046-b-transfer}(1), with
\(L=\{<\}\), \(\kappa=\lambda=\omega\) and \(\mathcal Q=\{Q_{\mathrm I}\}\),
the sentence \(\mathrm{WO}\) of
\cref{thm:p03-w03-p03-046-b-wellorder}(1) is equivalent on every infinite
\(L\)-structure to a sentence \(\mathrm{WO}'\) of
\((\mathrm{Th}(<),Q_{\mathrm I})\). Then \(\mathrm{WO}'\) holds in exactly
the infinite well-orders among the infinite linear orders, contrary to
\cref{thm:p03-w03-p03-046-b-wellorder}(2).

\emph{\(Q_{1-1}\not\le Q_{\mathrm{mon}}\).} Suppose
\(Q_{1-1}\le Q_{\mathrm{mon}}\). Let \(\Sigma''\) be the sentence of
\cref{thm:p03-w03-p03-046-b-uncountable}, whose infinite models are the
uncountable sets, and \(\Sigma''^\sharp\) the
\((\mathrm{Th}(=),Q_{\mathrm{mon}})\)-sentence of
\cref{thm:p03-w03-p03-046-b-relativize}. Since \(\Sigma''\) has an infinite
model, so does \(\Sigma''^\sharp\), and by
\cref{thm:p03-w03-p03-046-b-mso-sets}(2) \(\Sigma''^\sharp\) then holds in
every infinite set, in particular in \(\omega\). So some guard parameters
in \(\omega\) give a set \(\theta_0(\omega,\bar a_0,\bar S_0)\) that is a
model of \(\Sigma''\); this set is infinite by clause~(i) and countable as a
subset of \(\omega\), which contradicts the choice of \(\Sigma''\).
\end{proof}
\end{annexitem}

\begin{annexitem}{annex:p03-w03-p03-046-b-turing}{thm:p03-w03-p03-046-b-turing}{Recursive syntactic interpretations and Turing degrees}
\begin{proof}
Fix G\"odel numberings under which the codes of \(T_1\)-sentences and of
\(T_2\)-sentences form recursive sets, and let \(X_i\) be the set of codes
of \(T_i\)-sentences \(\phi\) with \(T_i\vdash\phi\). Let \(h\) be the
recursive function sending the code of \(\phi\) to the code of \(\phi^*\).
The \(T_1\)-sentence \(\phi_0:=\exists x\,x\neq x\) is false in every
model, and \(T_1\) has a model, so \(T_1\nvdash\phi_0\) and hence
\(T_2\nvdash\phi_0^*\). Put \(h'(e)=h(e)\) if \(e\) codes a
\(T_1\)-sentence and \(h'(e)=h(\ulcorner\phi_0\urcorner)\) otherwise; \(h'\)
is recursive. If \(e\) codes \(\phi\), then \(e\in X_1\) iff
\(T_1\vdash\phi\) iff \(T_2\vdash\phi^*\) iff \(h'(e)\in X_2\), by
\cref{def:p03-w03-p03-046-b-syntactic}; if \(e\) codes no \(T_1\)-sentence,
then \(e\notin X_1\) and \(h'(e)\notin X_2\). So \(h'\) is a many-one
reduction of \(X_1\) to \(X_2\); a many-one reduction is a Turing
reduction (ask the oracle whether \(h'(e)\in X_2\)), so the Turing degree
of \(X_1\) is at most that of \(X_2\).
\end{proof}
\end{annexitem}

\begin{annexitem}{annex:p03-w03-p03-046-b-codable}{thm:p03-w03-p03-046-b-codable}{Codable theories (Theorem~1.3.2)}
\begin{proof}
Let \(M\entails T\), \(n\ge2\), infinite \(B_0,\dots,B_{n-1},C\subseteq M\)
and \(\phi(x,y_0,\dots,y_{n-1},\bar d)\) be as in
\cref{def:p03-w03-p03-046-b-codable}. Fix \(b_j\in B_j\) for \(2\le j<n\)
and put \(F(u_0,u_1,x):=\phi(x,u_0,u_1,b_2,\dots,b_{n-1},\bar d)\) and
\(C'=\{x:\exists(u_0,u_1)\in B_0\times B_1\ F(u_0,u_1,x)\}\subseteq C\). For
each \((u_0,u_1)\in B_0\times B_1\) there is exactly one \(x\) with
\(F(u_0,u_1,x)\), it lies in \(C'\), distinct pairs give distinct \(x\)
(they extend to distinct \(n\)-tuples), and \(C'\) is the set of values;
\(C'\) is infinite because \(B_0\times B_1\) is. So \(T\) admits coding in
the sense of \cref{def:p03-w03-p03-046-pairing}(2), with the sets
\(B_0,B_1,C'\), the formula \(F\) (in which the unary predicates do not
occur) and the constants \(b_2,\dots,b_{n-1},\bar d\).

\Cref{thm:p03-w03-p03-046-coding} then gives
\((T_\infty,\mathrm{2nd})\le(T,\mathit{Mon})\) in the sense of
\cref{def:p03-w03-p03-046-interpretation}. That lemma assumes that \(T\) has
no finite models, but its proof does not use this: Step~1 only needs a
model of \(T\) with the infinite sets, and in clause~(1) the domain is
infinite because the guard codes a one-to-one self-map of it that is not
onto. Under the correspondence stated after
\cref{def:p03-w03-p03-046-b-reading} (guard \(\pi\), domain formula
\(\theta\), coding formula \(\chi_0\)), and since \(T_\infty\) has no
relation symbols to interpret and the parameters of \(\mathit{Mon}\) and
\(Q_{\mathrm{mon}}\) are the same (all subsets), this is
\(Q_{\mathrm{II}}\le(T,Q_{\mathrm{mon}})\) in the sense of
\cref{def:p03-w03-p03-046-b-interpretation}.
\end{proof}
\end{annexitem}

\begin{annexitem}{annex:p03-w03-p03-046-b-equiv}{thm:p03-w03-p03-046-b-equiv}{Equivalence relations with infinitely many infinite classes (Lemma~1.3.3)}
\begin{proof}
\Cref{thm:p03-w03-p03-046-equiv} gives an interpretation
\((\pi,\theta,\chi_0)\) of \((T_\infty,\mathrm{2nd})\) in \((T,\text{1-1})\),
whose relation parameters are involutions. Its proof does not use the
absence of finite models of \(T\): the compactness step only needs the
given model, and the domain of the interpretation is infinite because the
guard of \cref{thm:p03-w03-p03-046-equiv-uniform} codes a one-to-one
self-map of it that is not onto.

Replace involution parameters by pairing parameters. The formula
\(\theta_{\mathrm{inv}}(x,y;\bar p)\) of the proof of
\(Q_{\mathrm{inv}}\le Q_{1-1}\) in
\cref{thm:p03-w03-p03-046-b-permutations} has no guard, and on every
infinite set its values, as \(\bar p\) ranges over points and pairings,
are exactly the involutions. In \(\pi\), \(\theta\) and \(\chi_0\) replace
each involution variable \(\sigma\) by \(\theta_{\mathrm{inv}}(\cdot,\cdot;\bar
p_\sigma)\) with a fresh tuple \(\bar p_\sigma\), giving \(\pi'\),
\(\theta'\), \(\chi_0'\). If \(M\entails T\) and pairing parameters satisfy
\(\pi'\), then by (F2) the involutions \(\sigma=\theta_{\mathrm{inv}}(M;\bar
p_\sigma)\) satisfy \(\pi\); so \(\theta'(M)=\theta(M,\bar\sigma)\) is
infinite, \(M\) is infinite, and as the coding tuples \(\bar p'\) range
over all points and pairings the involutions they define range over all
involutions of \(M\), so the relations \(\chi_0'(M,\dots)\) are exactly
the binary relations on \(\theta'(M)\). For each \(\kappa\), the model and
involutions given by the original clause~(2) are reached by suitable
pairing tuples. This is \(Q_{\mathrm{II}}\le(T,Q_{1-1})\) in the sense of
\cref{def:p03-w03-p03-046-b-interpretation}.
\end{proof}
\end{annexitem}
